\ifreviewchanges
\clearpage
\chapter*{Review Question Source Register}
\addcontentsline{toc}{chapter}{Review Question Source Register}
\footnotesize
\renewcommand{\arraystretch}{1.12}
This review-only register assigns a stable \textbf{Q} ID to each live question row
in the reviewed paper.  The ID appears before the question in the tracked-change
PDF, and the third column gives the source question or comment text used to ground
the row.  The complete original source inventory is listed separately as
\textbf{OQ001}--\textbf{OQ158} in Table~\ref{tab:original-source-question-register},
and the original review comments are listed as \textbf{RC001}--\textbf{RC060} in
Table~\ref{tab:original-review-comment-register}.  The live \textbf{Q} IDs are
intentionally non-contiguous: resolved setup questions, including the former
boundary-attachment, field-definition, residual-definition, and RH-conversion rows,
were dissolved into the main prose and mapping tables rather than renumbered.\par\medskip
\begin{longtable}{@{}L{0.10\textwidth} L{0.30\textwidth} L{0.52\textwidth}@{}}
\caption{Live review question IDs and original source question/comment text.}\label{tab:review-question-source-register}\\
\toprule
ID & Paper question & Original source question/comment text\\
\midrule
\endfirsthead
\caption[]{Live review question IDs and original source question/comment text (continued).}\\
\toprule
ID & Paper question & Original source question/comment text\\
\midrule
\endhead
\bottomrule
\endlastfoot
\textbf{Q006} & Which timestamp and time-origin convention produced the active niche curve \(p_E(t)\)? & Source inventory 1.3: What is the current source of the active niche boundary pressure curve: XML only, provider spreadsheet, RH-derived Kelvin reconstruction, fitted seasonal curve, pressure sensor data, or another processed measurement? Source inventory 1.4: How should model-output comparisons handle time: OGS output timestep, calendar date, measurement timestamp, survey epoch, monthly ERT schedule, seasonal campaign, interpolation, extrapolation, and early numerical adjustment from inconsistent initial states? User revision: change this to a query to check timestamps, because initial experiments show that at least ERT suggests it is shifted in time.\\
\addlinespace[0.22em]
\textbf{Q012} & Does the initial pressure encode post-excavation drainage near the niche? & Source inventory 1.3: For hydraulic initial state, should the initial pressure be uniform 1.5 MPa, a heterogeneous field, or a post-excavation field already drained near the niche by previous excavations and boreholes? Annotated PDF A06: chceme si vyjasnit.\\
\addlinespace[0.22em]
\textbf{Q013} & Does the model simulate staged excavation or gradual unloading? & Source inventory 1.3: Is the niche/tunnel excavation represented as a staged process, or does the model start with an already-open geometry and prescribed initial fields? Annotated PDF A05: kolik casu je od exkavace?\\
\addlinespace[0.22em]
\textbf{Q014} & Which prestress magnitude, angle, and excavation reference should be used if prestress is activated? & Source inventory 2.5: What in-situ stress tensor and stress sign/reference convention should be used for a deep tunnel/niche model, and are initial\_stress, ic\_sigma0, and load\_top active or only inactive provenance? User follow-up: There is usually an angle in the prestress. User revision: steer this to prestress query and merge with Q15. Annotated PDF A07/A08: chceme vyjasnit.\\
\addlinespace[0.22em]
\textbf{Q016} & Which early-time model-data mismatch is a physical signal rather than initialization artefact? & Source inventory 1.4: How should model-output comparisons handle time: OGS output timestep, calendar date, measurement timestamp, survey epoch, monthly ERT schedule, seasonal campaign, interpolation, extrapolation, and early numerical adjustment from inconsistent initial states? User revision: reformulate to be more understandable and tied to the text. Annotated PDF A09: chceme vyjasnit.\\
\addlinespace[0.22em]
\textbf{Q020} & Which governing terms are explicitly inactive, and where do they appear in the equations? & Source inventory 1.1: Which terms are absent or inactive even if they appear in related TRM theory: vapor diffusion, hydrostatic gravity, active top mechanical load, active initial-stress/prestress parameter, and the outside pressure parameter that is defined but not attached? User revision: add specific equation numbers this addresses. Annotated PDF A13: ok, pokud nechceme menit model.\\
\addlinespace[0.22em]
\textbf{Q021} & Which parameter family is allowed to vary in the first workflow? & Source inventory 1.2: Which parameters are fixed for the current workflow, which are active/released, which are inactive/commented provenance, and which are later-stage candidates gated by measurement semantics? Earlier note: "Which parameter family is released in the first workflow?" nerozumim tomu, co znamena released.\\
\addlinespace[0.22em]
\textbf{Q022} & How do candidate-run parameter changes enter OGS? & Source inventory 1.2: If a parameter changes between candidate runs, is it an OGS model-input field or XML replacement generated externally, not a silent change to governing equations?\\
\addlinespace[0.22em]
\textbf{Q025} & Are mechanical parameters fitted now? & Source inventory 2.5: Relation to OGS: Are orthotropic elasticity, Biot coefficient, Bishop b(S\_l)=S\_l, and saturation-dependent swelling fixed model inputs in the current workflow? Source inventory 2.5: Is there enough displacement, crack, levelling, extensometer, pressure, or laser-scan data to release mechanical parameters, or should they remain fixed/context until numeric HM residuals exist?\\
\addlinespace[0.22em]
\textbf{Q026} & Are thermal or liquid properties fitted now? & Source inventory 2.6: Relation to OGS: Which thermal/fluid fields are XML constants or phase-property laws, and which are not defensible inversion parameters now? Source inventory 2.6: Is the active CTE value physically meaningful, inactive, typo/copied heat capacity, or a provenance issue that must be confirmed before interpreting thermal-mechanical coupling?\\
\addlinespace[0.22em]
\textbf{Q027} & Should any material parameter be time-dependent in this first pass? & Source inventory 2.7: Which material inputs may need time dependence over the multi-year experiment, especially EDZ permeability decreasing through clay healing/sealing, fracture closure/opening, porosity/retention changes, or boundary-curve changes? Source inventory 2.7: What evidence would justify releasing time-dependent material fields: repeated same-support permeability, NMR/Taupe/ERT trend drift under fixed fields, crack-aperture time series, Geoscope/laser/levelling data, pressure-response changes, or geochemical/mineralogical sealing evidence?\\
\addlinespace[0.22em]
\textbf{Q032} & How are output times tied to observations or boundary curves? & Source inventory 1.4: How should model-output comparisons handle time: OGS output timestep, calendar date, measurement timestamp, survey epoch, monthly ERT schedule, seasonal campaign, interpolation, extrapolation, and early numerical adjustment from inconsistent initial states? User revision: check timestamps because initial experiments suggest at least ERT is shifted in time.\\
\addlinespace[0.22em]
\textbf{Q042} & Does the same bedding angle control mechanics, permeability, and electrical anisotropy? & Source inventory 2.2 and 2.5: Does the same bedding angle control Darcy permeability anisotropy, ERT electrical anisotropy, and mechanical orthotropy, or are these separate modelling choices? Annotated PDF A21: potrebujeme vedet, plus prestress. User follow-up: There is usually an angle in the prestress.\\
\addlinespace[0.22em]
\textbf{Q044} & How should faults, cracks, and EDZ features be classified here? & Source inventory 2.7: Should faults/fractures/cracks become hard geometry inputs, structural priors, EDZ/fault material masks, permeability/porosity fields, mechanical caveats, or qualitative diagnostics? Annotated PDF A24: pozdejsi faze nebo od zacatku?\\
\addlinespace[0.22em]
\textbf{Q045} & How should tests that rapidly oscillate between the lowest and highest interpreted permeability be understood? & Source inventory 3.2: Time of measurement: Which campaign/workbook date applies to each interpreted row and raw pressure-decay row, and should repeated support rows be treated as independent, duplicates, or time evolution? Source inventory 3.2: Uncertainty/use: Should the likelihood remain rowwise Gaussian, robust, support-cell aggregated, capped, or duplicate-weighted, given strong same-support conflicts? User revision: drop current Q45--Q49 and replace with one question about tests that very quickly oscillate between lowest and highest and how reliable this can be.\\
\addlinespace[0.22em]
\textbf{Q050} & What exactly does the supplied NMR value measure? & Source inventory 3.3: Meaning: Does NMR measure total NMR-visible water, mobile/free water, bound/interlayer water, or another hydrogen/water proxy in the supplied CD-A data? Annotated PDF A27: !!!\\
\addlinespace[0.22em]
\textbf{Q051} & How should NMR be linked to OGS output now? & Source inventory 3.3: OGS tie: Should comparison be to OGS saturation S\_l, theta = n*S\_l, corrected/free-water theta, label-bias corrected theta, or within-label trend/anomaly only? Annotated PDF A28: !!!\\
\addlinespace[0.22em]
\textbf{Q052} & Can bound/interlayer water be anchored simply? & Source inventory 3.3: Bound water: How should bound/interlayer water not active in OGS saturation be filtered, modelled, estimated, held constant, or represented as bias/uncertainty? Annotated PDF A29: asi spise pro nas?\\
\addlinespace[0.22em]
\textbf{Q053} & Could bound water be post-processed from OGS saturation? & Source inventory 3.3: Bound water: How should bound/interlayer water not active in OGS saturation be filtered, modelled, estimated, held constant, or represented as bias/uncertainty? Annotated PDF A30: asi spise pro nas?\\
\addlinespace[0.22em]
\textbf{Q054} & What would fully ground NMR residuals? & Source inventory 3.3: Uncertainty/use: What uncertainty floor and confidence-interval conversion should be used, and how should rows above fixed porosity be interpreted? Annotated PDF A31: asi spise pro nas.\\
\addlinespace[0.22em]
\textbf{Q055} & What ERT value should be compared? & Source inventory 3.4: Meaning: Is the primary ERT observable resistivity, log resistivity, resistivity change, phase angle, or a derived water-content field?\\
\addlinespace[0.22em]
\textbf{Q056} & How does ERT link to OGS now? & Source inventory 3.4: OGS tie: Should OGS saturation or theta = n*S\_l be converted to resistivity through an Archie-type, empirical CD-A-specific, NMR-calibrated, or clay-specific relation?\\
\addlinespace[0.22em]
\textbf{Q057} & Does bound water conduct like free water? & Source inventory 3.4: Bound water and clay conduction: Does bound/locked water conduct current like free pore water, and how do clay surface conduction, salinity, porosity, pore-water conductivity, cracks/faults, and saturation history affect resistivity? Annotated PDF A32: !!!\\
\addlinespace[0.22em]
\textbf{Q058} & Does electrical anisotropy matter? & Source inventory 3.4: Anisotropy: Does electrical anisotropy matter, and should it connect to bedding, fractures/faults, or permeability anisotropy? Annotated PDF A33: !!!\\
\addlinespace[0.22em]
\textbf{Q059} & What would fully ground ERT likelihood use? & Source inventory 3.4: Uncertainty/use: What uncertainty/covariance/aggregation model prevents over-weighting dense correlated ERT cells, and should ERT remain diagnostic until transform/support/uncertainty are approved?\\
\addlinespace[0.22em]
\textbf{Q086} & Can we obtain raw ERT data for all campaigns, or at least the first two years? & User revision: add a Q into ERT, a query to obtain raw data for all measurements or at least the starting two years for quicker tests. Source inventory 3.4: Which ERT timesteps, monthly campaigns, folders, and VTK fields correspond to OGS output times? Source inventory 3.4: What uncertainty/covariance/aggregation model prevents over-weighting dense correlated ERT cells?\\
\addlinespace[0.22em]
\textbf{Q060} & What does the Taupe workbook value represent? & Source inventory 3.5: Meaning: Does TAUPE mean the Taupe/TDR workbook stream in the repo, and what physical quantity does it represent? Annotated PDF A34: !!!\\
\addlinespace[0.22em]
\textbf{Q061} & How should Taupe/TDR be linked to OGS now? & Source inventory 3.5: OGS tie: Should Taupe/TDR compare to absolute theta = n*S\_l, saturation S\_l, band-averaged theta anomaly, or grouped trend diagnostics only?\\
\addlinespace[0.22em]
\textbf{Q062} & When can Taupe/TDR become a calibrated likelihood residual? & Earlier text note and source inventory 3.5: "Can it become an absolute residual?" nerozumim pojmu "absolute residual". Units: What are the workbook units, calibration equations, constants, baseline normalization, and sensor/band-specific scaling rules? User revision: reformulate to convey what residual means, i.e. state in terms of measurement and likelihood term to inform posterior. Annotated PDF A35: nerozumim.\\
\addlinespace[0.22em]
\textbf{Q063} & How should bands and sensors be weighted? & Source inventory 3.5: Uncertainty/use: What uncertainty and weighting should be assigned by sensor, band, time, calibration quality, support projection, and grouping? Annotated PDF A36: nerozumim.\\
\addlinespace[0.22em]
\textbf{Q064} & What would fully ground Taupe/TDR use? & Source inventory 3.5: Units: What are the workbook units, calibration equations, constants, baseline normalization, and sensor/band-specific scaling rules?\\
\addlinespace[0.22em]
\textbf{Q067} & Which sensors define the open-niche boundary? & Source inventory 3.6: Position in domain: Which RH/suction sensors correspond to the open niche or open twin, and which should inform the niche boundary rather than interior cell residuals? Annotated PDF A40: spise bych se ptal na nejpouzitelnejsi sensor.\\
\addlinespace[0.22em]
\textbf{Q069} & What would fully ground the pressure boundary? & Source inventory 3.6: Curve provenance: Why does the local RH-derived envelope differ from the active seasonal pressure curve, and what source table/script/sensor-screening policy generated the active curve? Uncertainty/use: What uncertainty should be assigned if RH becomes a boundary or retention likelihood, and how do we avoid double-counting the same data as both forcing and validation?\\
\addlinespace[0.22em]
\textbf{Q070} & Do direct pore-pressure measurements exist for this model? & Source inventory 3.7: Meaning: Do direct pressure, hydraulic-head, pore-pressure, or mini-piezometer measurements exist, including the two more distant sensors remembered from Gesa material?\\
\addlinespace[0.22em]
\textbf{Q071} & How should pressure be linked to OGS? & Source inventory 3.7: OGS tie: Can these measurements compare directly to OGS liquid pressure p\_l, or do they require head/elevation/gauge/reference conversion?\\
\addlinespace[0.22em]
\textbf{Q072} & Can pressure data constrain initial and boundary conditions? & Source inventory 3.7: Uncertainty/use: Are they reliable enough for hard residuals, or only qualitative validation because of distance, noise, sensor failures, maintenance, calibration, or 2D-model error?\\
\addlinespace[0.22em]
\textbf{Q073} & What 2D-model caveat is needed? & Source inventory 3.7: Position in domain: Where are the sensors in coordinates, elevation, borehole/support geometry, distance from the niche, and 2D model frame? Annotated PDF A41: prakticky chceme vedet, jak interpretovat data(2D/3D)?\\
\addlinespace[0.22em]
\textbf{Q074} & What would fully ground pressure residuals? & Source inventory 3.7: What/how measured: What pressure quantity is measured: pore pressure, hydraulic head, gauge pressure, absolute pressure, suction/capillary pressure, or a processed Geoscope value? Units: Are values Pa, MPa, bar, head in m, gauge pressure, absolute pressure, or relative to atmospheric/gas reference?\\
\addlinespace[0.22em]
\textbf{Q075} & What does the large-fault source measure? & Source inventory 3.8: Meaning: Does the faulty/.dat source mean large fault/fault-surface/fracture geometry rather than niche-surface crackmeter data?\\
\addlinespace[0.22em]
\textbf{Q076} & How should faults link to OGS now? & Source inventory 3.8: OGS tie: Should the projected fault become a geometry input, structural prior, EDZ/fault material class, permeability or porosity mask, retention/storage prior, mechanical caveat, or qualitative context? Annotated PDF A42: Asi chceme jen navrhnout, ze budeme pracovat s trhlinami jako extra nahodnym polem/domenou a nechat si to schvalit?\\
\addlinespace[0.22em]
\textbf{Q077} & Can a fault modify permeability or porosity? & Source inventory 3.8: What/how measured: Does the source describe geometry, aperture, hydraulic aperture, mechanical aperture, fracture opening, displacement, fault-zone class, or only visualization geometry?\\
\addlinespace[0.22em]
\textbf{Q078} & How should 3D geometry be reduced to the 2D slice? & Source inventory 3.8: Projection/use: What 3D-to-2D projection is accepted, and how do we avoid overfitting a 2D permeability field to unresolved 3D structures?\\
\addlinespace[0.22em]
\textbf{Q079} & What would fully ground quantitative fault use? & Source inventory 3.8: Position in domain: Which exact .dat file/source contains the 3D fault geometry, coordinate system, units, axes, trace/surface/polygon/point-cloud representation, and relation to the current Niche 4 2D slice?\\
\addlinespace[0.22em]
\textbf{Q080} & What does each HM stream measure? & Source inventory 3.9: Meaning: What visible or measurable crack/opening/deformation data exist on the niche surface, separate from large 3D fault geometry? Annotated PDF A43: nemam predstavu, jak tahla data pouzit.\\
\addlinespace[0.22em]
\textbf{Q081} & How should crackmeters link to OGS? & Source inventory 3.9: OGS tie: Should comparison use OGS displacement, strain, stress, derived crack opening from displacement difference, pressure response, swelling/deformation pattern, or only mechanical plausibility screening?\\
\addlinespace[0.22em]
\textbf{Q082} & Can visible cracks justify permeability changes? & Source inventory 3.9: Material-field tie: Can measured cracks/openings justify local permeability anomaly, porosity/damage zone, hydraulic aperture relation, fracture/EDZ mask, mechanical stiffness reduction, or time-dependent permeability/porosity change?\\
\addlinespace[0.22em]
\textbf{Q083} & How should levelling, extensometers, and laser scans be used now? & Source inventory 3.9: What/how measured: Are quantities crack width, aperture, crack opening displacement, displacement jump, strain, surface displacement, laser-scan difference, convergence, levelling vertical displacement, pressure response, or qualitative crack class?\\
\addlinespace[0.22em]
\textbf{Q084} & Which other measurements can enter the OGS workflow? & Source inventory 3.10: Other measurements: What additional measurement streams exist in the data beyond those listed here, and for each one what are where, when, units, measured quantity, and OGS input/output/prior/diagnostic status?\\
\addlinespace[0.22em]
\textbf{Q085} & What would fully ground HM residuals? & Source inventory 3.9 and 3.10: Uncertainty/use: What numeric exports, support geometry, reference-zero/sign convention, quality flags, registration uncertainty, sensor noise, and 2D projection error are needed before hard residuals? Activation/use: Which streams are active likelihood terms now, which are diagnostic, which are boundary/input provenance, which are support/prior only, and which are blocked by missing numeric exports or metadata?\\
\addlinespace[0.22em]
\end{longtable}
% BEGIN GENERATED DETAILED MEASUREMENT QUESTION REGISTER
\clearpage
\section*{Detailed Measurement Question Source Register}
\addcontentsline{toc}{section}{Detailed Measurement Question Source Register}
This review-only register assigns a stable \textbf{DQ} ID to every detailed measurement-question row in the measurement-question inventory.  The full source-question text is repeated here so each detailed row can be discussed by ID without searching the original inventory.\par\medskip
\scriptsize
\renewcommand{\arraystretch}{1.10}
\begin{longtable}{@{}L{0.09\textwidth} L{0.31\textwidth} L{0.52\textwidth}@{}}
\caption{Detailed measurement-question IDs and original source question/comment text.}\label{tab:detailed-measurement-question-source-register}\\
\toprule
ID & Paper question & Original source question/comment text\\
\midrule
\endfirsthead
\caption[]{Detailed measurement-question IDs and original source question/comment text (continued).}\\
\toprule
ID & Paper question & Original source question/comment text\\
\midrule
\endhead
\bottomrule
\endlastfoot
\textbf{DQ001} & Which generated reduction defines the measurement positions? & Source inventory 3.1: For each measurement, where is it taken: sensor id, coordinates, borehole/support geometry, niche side, depth/elevation, interval/band/surface, and relation to the current 2D OGS slice? Source inventory 3.1: For each measurement, what preprocessing is required: unit conversion, coordinate transform, 3D-to-2D projection, support averaging, baseline/reference-zero definition, time interpolation, filtering, and uncertainty/correlation model?\\
\addlinespace[0.18em]
\textbf{DQ002} & Which niche is represented in the current note? & User follow-up and source inventory 3.1/3.10: User follow-up: is the open niche the one at 0 or the one around 12? which one we are approximating, also which sensor data are we considering int he note pdf? we should only consider the measurements around the niche with the ERT sensor data, (open one) User follow-up and source inventory 3.1/3.10: For each measurement, where is it taken: sensor id, coordinates, borehole/support geometry, niche side, depth/elevation, interval/band/surface, and relation to the current 2D OGS slice? User follow-up and source inventory 3.1/3.10: Position in domain: Which coordinate rows, borehole endpoints, support points, line samples, ERT mesh cells, NMR/Taupe supports, and fault/bedding geometries are inside, outside, or near the current 2D mesh?\\
\addlinespace[0.18em]
\textbf{DQ003} & Can every measurement be treated as a point in the OGS mesh? & Source inventory 3.1: For each measurement, where is it taken: sensor id, coordinates, borehole/support geometry, niche side, depth/elevation, interval/band/surface, and relation to the current 2D OGS slice? Source inventory 3.1: For each measurement, what preprocessing is required: unit conversion, coordinate transform, 3D-to-2D projection, support averaging, baseline/reference-zero definition, time interpolation, filtering, and uncertainty/correlation model?\\
\addlinespace[0.18em]
\textbf{DQ004} & Which streams have available measured-data plots? & Source inventory 3.1/3.10: For each measurement, what exactly does it measure before modelling interpretation, and by what physical or processing method? Source inventory 3.1/3.10: For each measurement, what are the units, reference convention, sign convention, calibration state, raw/processed status, source files, and quality flags? Source inventory 3.1/3.10: Other measurements: What additional measurement streams exist in the data beyond those listed here, and for each one what are where, when, units, measured quantity, and OGS input/output/prior/diagnostic status?\\
\addlinespace[0.18em]
\textbf{DQ005} & What raw source file defines this stream? & Source inventory 3.1/3.10: For each measurement, what exactly does it measure before modelling interpretation, and by what physical or processing method? Source inventory 3.1/3.10: For each measurement, what are the units, reference convention, sign convention, calibration state, raw/processed status, source files, and quality flags? Source inventory 3.1/3.10: Other measurements: What additional measurement streams exist in the data beyond those listed here, and for each one what are where, when, units, measured quantity, and OGS input/output/prior/diagnostic status?\\
\addlinespace[0.18em]
\textbf{DQ006} & What processing step produced the table used here? & Source inventory 3.1/3.10: For each measurement, what exactly does it measure before modelling interpretation, and by what physical or processing method? Source inventory 3.1/3.10: For each measurement, what are the units, reference convention, sign convention, calibration state, raw/processed status, source files, and quality flags? Source inventory 3.1/3.10: Other measurements: What additional measurement streams exist in the data beyond those listed here, and for each one what are where, when, units, measured quantity, and OGS input/output/prior/diagnostic status?\\
\addlinespace[0.18em]
\textbf{DQ007} & What exact spatial support does one row represent? & Source inventory 3.1: For each measurement, where is it taken: sensor id, coordinates, borehole/support geometry, niche side, depth/elevation, interval/band/surface, and relation to the current 2D OGS slice? Source inventory 3.1: For each measurement, what preprocessing is required: unit conversion, coordinate transform, 3D-to-2D projection, support averaging, baseline/reference-zero definition, time interpolation, filtering, and uncertainty/correlation model?\\
\addlinespace[0.18em]
\textbf{DQ008} & What timestamp or survey epoch does one row represent? & Source inventory 3.1: For each measurement, when is it taken: timestamp, survey epoch, campaign date, averaging window, time zone/date convention, and relation to OGS time zero/output timesteps/boundary-curve times?\\
\addlinespace[0.18em]
\textbf{DQ009} & What unit and reference convention does the stored value use? & Source inventory 3.1: For each measurement, what are the units, reference convention, sign convention, calibration state, raw/processed status, source files, and quality flags?\\
\addlinespace[0.18em]
\textbf{DQ010} & What uncertainty or covariance should be assigned to the row? & Source inventory 3.1: For each measurement, what preprocessing is required: unit conversion, coordinate transform, 3D-to-2D projection, support averaging, baseline/reference-zero definition, time interpolation, filtering, and uncertainty/correlation model?\\
\addlinespace[0.18em]
\textbf{DQ011} & What OGS object does the stream ground? & Source inventory 3.1/3.10: For each measurement, how does it tie to OGS: input/boundary forcing, initial-condition constraint, material-field prior, model-output residual, diagnostic screen, qualitative caveat, or coordinate/support operator? Source inventory 3.1/3.10: Activation/use: Which streams are active likelihood terms now, which are diagnostic, which are boundary/input provenance, which are support/prior only, and which are blocked by missing numeric exports or metadata?\\
\addlinespace[0.18em]
\textbf{DQ012} & What must a row contain before an inverse run may use it? & Source inventory 3.1/3.10: For each measurement, how does it tie to OGS: input/boundary forcing, initial-condition constraint, material-field prior, model-output residual, diagnostic screen, qualitative caveat, or coordinate/support operator? Source inventory 3.1/3.10: Activation/use: Which streams are active likelihood terms now, which are diagnostic, which are boundary/input provenance, which are support/prior only, and which are blocked by missing numeric exports or metadata?\\
\addlinespace[0.18em]
\textbf{DQ013} & How should incomplete rows enter candidate comparisons? & Source inventory 3.1/3.10: For each measurement, how does it tie to OGS: input/boundary forcing, initial-condition constraint, material-field prior, model-output residual, diagnostic screen, qualitative caveat, or coordinate/support operator? Source inventory 3.1/3.10: Activation/use: Which streams are active likelihood terms now, which are diagnostic, which are boundary/input provenance, which are support/prior only, and which are blocked by missing numeric exports or metadata?\\
\addlinespace[0.18em]
\textbf{DQ014} & What inverse-problem failure does support uncertainty create? & Source inventory 3.1: For each measurement, where is it taken: sensor id, coordinates, borehole/support geometry, niche side, depth/elevation, interval/band/surface, and relation to the current 2D OGS slice? Source inventory 3.1: For each measurement, what preprocessing is required: unit conversion, coordinate transform, 3D-to-2D projection, support averaging, baseline/reference-zero definition, time interpolation, filtering, and uncertainty/correlation model?\\
\addlinespace[0.18em]
\textbf{DQ015} & What does the permeability stream measure? & Source inventory 3.2: Meaning: Are pulse tests interpreted as noisy scalar interval observations of intrinsic permeability/transmissibility, not as direct tensor components, hydraulic conductivity, relative permeability, or saturation? Source inventory 3.2: Uncertainty/use: Should the likelihood remain rowwise Gaussian, robust, support-cell aggregated, capped, or duplicate-weighted, given strong same-support conflicts?\\
\addlinespace[0.18em]
\textbf{DQ016} & Where and when is it used in the current OGS objective? & Source inventory 3.2: Meaning: Are pulse tests interpreted as noisy scalar interval observations of intrinsic permeability/transmissibility, not as direct tensor components, hydraulic conductivity, relative permeability, or saturation? Source inventory 3.2: Uncertainty/use: Should the likelihood remain rowwise Gaussian, robust, support-cell aggregated, capped, or duplicate-weighted, given strong same-support conflicts?\\
\addlinespace[0.18em]
\textbf{DQ017} & How is the permeability residual formed? & Source inventory 3.2: OGS tie: Should each row compare to interval/support average of e\textasciicircum{}T K e from the OGS intrinsic permeability tensor field k\_i\_rd, with explicit support weights and uncertainty? Source inventory 3.2: Uncertainty/use: Should the likelihood remain rowwise Gaussian, robust, support-cell aggregated, capped, or duplicate-weighted, given strong same-support conflicts?\\
\addlinespace[0.18em]
\textbf{DQ018} & Is each permeability value intrinsic liquid-equivalent permeability? & Source inventory 3.2: Meaning: Are pulse tests interpreted as noisy scalar interval observations of intrinsic permeability/transmissibility, not as direct tensor components, hydraulic conductivity, relative permeability, or saturation? Source inventory 3.2: Units: Are values intrinsic permeability in m2, transmissibility, log10 permeability, hydraulic conductivity, or workbook-specific units, and how are zeros/missing values handled?\\
\addlinespace[0.18em]
\textbf{DQ019} & Was a Klinkenberg or gas-slip correction applied? & Source inventory 3.2: What/how measured: What gas, pressure-decay protocol, correction, Klinkenberg/slip treatment, interval length, and interpretation method produced each permeability value?\\
\addlinespace[0.18em]
\textbf{DQ020} & What gas type and mean pressure were used for each test? & Source inventory 3.2: What/how measured: What gas, pressure-decay protocol, correction, Klinkenberg/slip treatment, interval length, and interpretation method produced each permeability value?\\
\addlinespace[0.18em]
\textbf{DQ021} & What temperature assumption was used in the gas-test interpretation? & Source inventory 3.2: What/how measured: What gas, pressure-decay protocol, correction, Klinkenberg/slip treatment, interval length, and interpretation method produced each permeability value?\\
\addlinespace[0.18em]
\textbf{DQ022} & Which pressure-decay fits are failed or low-signal? & Source inventory 3.2: What/how measured: What gas, pressure-decay protocol, correction, Klinkenberg/slip treatment, interval length, and interpretation method produced each permeability value?\\
\addlinespace[0.18em]
\textbf{DQ023} & Which rows may be affected by leakage? & Source inventory 3.2: What/how measured: What gas, pressure-decay protocol, correction, Klinkenberg/slip treatment, interval length, and interpretation method produced each permeability value?\\
\addlinespace[0.18em]
\textbf{DQ024} & What interval centre and interval length belong to each row? & Source inventory 3.2: Position in domain: What are the exact borehole intervals, endpoints, orientations, active mesh cells, and blocked historical endpoint geometries for BCD-A24/25/26/27, BCD-A32/A33, BFM-D19, and related rows? Source inventory 3.2: OGS tie: Should each row compare to interval/support average of e\textasciicircum{}T K e from the OGS intrinsic permeability tensor field k\_i\_rd, with explicit support weights and uncertainty?\\
\addlinespace[0.18em]
\textbf{DQ025} & What borehole trace and direction belong to each row? & Source inventory 3.2: Position in domain: What are the exact borehole intervals, endpoints, orientations, active mesh cells, and blocked historical endpoint geometries for BCD-A24/25/26/27, BCD-A32/A33, BFM-D19, and related rows? Source inventory 3.2: OGS tie: Should each row compare to interval/support average of e\textasciicircum{}T K e from the OGS intrinsic permeability tensor field k\_i\_rd, with explicit support weights and uncertainty?\\
\addlinespace[0.18em]
\textbf{DQ026} & What 3D volume is sampled by one pulse-test row? & Source inventory 3.2: Position in domain: What are the exact borehole intervals, endpoints, orientations, active mesh cells, and blocked historical endpoint geometries for BCD-A24/25/26/27, BCD-A32/A33, BFM-D19, and related rows? Source inventory 3.2: OGS tie: Should each row compare to interval/support average of e\textasciicircum{}T K e from the OGS intrinsic permeability tensor field k\_i\_rd, with explicit support weights and uncertainty?\\
\addlinespace[0.18em]
\textbf{DQ027} & Is the directional interval operator accepted for final scoring? & Source inventory 3.2: Position in domain: What are the exact borehole intervals, endpoints, orientations, active mesh cells, and blocked historical endpoint geometries for BCD-A24/25/26/27, BCD-A32/A33, BFM-D19, and related rows? Source inventory 3.2: OGS tie: Should each row compare to interval/support average of e\textasciicircum{}T K e from the OGS intrinsic permeability tensor field k\_i\_rd, with explicit support weights and uncertainty?\\
\addlinespace[0.18em]
\textbf{DQ028} & Which bedding angle is used for permeability anisotropy? & Source inventory 2.2: What anisotropy model is defensible: fixed bedding-informed angle, global anisotropy ratio, spatially variable tensor orientation, fracture/fault/EDZ masks, or scalar magnitude field only? Source inventory 2.2: Does the same bedding angle control Darcy permeability anisotropy, ERT electrical anisotropy, and mechanical orthotropy, or are these separate modelling choices?\\
\addlinespace[0.18em]
\textbf{DQ029} & Which repeated tests sample the same support? & Source inventory 3.2: Position in domain: What are the exact borehole intervals, endpoints, orientations, active mesh cells, and blocked historical endpoint geometries for BCD-A24/25/26/27, BCD-A32/A33, BFM-D19, and related rows? Source inventory 3.2: OGS tie: Should each row compare to interval/support average of e\textasciicircum{}T K e from the OGS intrinsic permeability tensor field k\_i\_rd, with explicit support weights and uncertainty?\\
\addlinespace[0.18em]
\textbf{DQ030} & How are several-log10 repeated-test disagreements classified? & Source inventory 3.2: Meaning: Are pulse tests interpreted as noisy scalar interval observations of intrinsic permeability/transmissibility, not as direct tensor components, hydraulic conductivity, relative permeability, or saturation? Source inventory 3.2: Uncertainty/use: Should the likelihood remain rowwise Gaussian, robust, support-cell aggregated, capped, or duplicate-weighted, given strong same-support conflicts?\\
\addlinespace[0.18em]
\textbf{DQ031} & Are 2021 and 2024 rows both material observations? & Source inventory 3.2: Time of measurement: Which campaign/workbook date applies to each interpreted row and raw pressure-decay row, and should repeated support rows be treated as independent, duplicates, or time evolution?\\
\addlinespace[0.18em]
\textbf{DQ032} & How are older rows without endpoint geometry handled? & Source inventory 3.2: Position in domain: What are the exact borehole intervals, endpoints, orientations, active mesh cells, and blocked historical endpoint geometries for BCD-A24/25/26/27, BCD-A32/A33, BFM-D19, and related rows? Source inventory 3.2: OGS tie: Should each row compare to interval/support average of e\textasciicircum{}T K e from the OGS intrinsic permeability tensor field k\_i\_rd, with explicit support weights and uncertainty?\\
\addlinespace[0.18em]
\textbf{DQ033} & Which unknowns does the permeability likelihood constrain directly? & Source inventory 3.2: OGS tie: Should each row compare to interval/support average of e\textasciicircum{}T K e from the OGS intrinsic permeability tensor field k\_i\_rd, with explicit support weights and uncertainty? Source inventory 3.2: Uncertainty/use: Should the likelihood remain rowwise Gaussian, robust, support-cell aggregated, capped, or duplicate-weighted, given strong same-support conflicts?\\
\addlinespace[0.18em]
\textbf{DQ034} & What can bias the fitted permeability field if left unresolved? & Source inventory 3.2: OGS tie: Should each row compare to interval/support average of e\textasciicircum{}T K e from the OGS intrinsic permeability tensor field k\_i\_rd, with explicit support weights and uncertainty? Source inventory 3.2: Uncertainty/use: Should the likelihood remain rowwise Gaussian, robust, support-cell aggregated, capped, or duplicate-weighted, given strong same-support conflicts?\\
\addlinespace[0.18em]
\textbf{DQ035} & What is the minimum final-use policy for this stream? & Source inventory 3.2: Meaning: Are pulse tests interpreted as noisy scalar interval observations of intrinsic permeability/transmissibility, not as direct tensor components, hydraulic conductivity, relative permeability, or saturation? Source inventory 3.2: Uncertainty/use: Should the likelihood remain rowwise Gaussian, robust, support-cell aggregated, capped, or duplicate-weighted, given strong same-support conflicts?\\
\addlinespace[0.18em]
\textbf{DQ036} & What physical quantity does NMR represent before modelling interpretation? & Source inventory 3.3: Meaning: Does NMR measure total NMR-visible water, mobile/free water, bound/interlayer water, or another hydrogen/water proxy in the supplied CD-A data? Source inventory 3.3: What/how measured: What NMR signal, T2 distribution, calibration, detuning caveat, campaign processing, or label/borehole conversion produces the reported water-content value?\\
\addlinespace[0.18em]
\textbf{DQ037} & What are the available time ranges and supports? & Source inventory 3.3: Position in domain: Which NMR labels, profiles, borehole points, wall/profile supports, and current Niche 4 mesh supports are active or outside support?\\
\addlinespace[0.18em]
\textbf{DQ038} & Why is raw absolute NMR water content not a final residual policy? & Source inventory 3.3: OGS tie: Should comparison be to OGS saturation S\_l, theta = n*S\_l, corrected/free-water theta, label-bias corrected theta, or within-label trend/anomaly only? Source inventory 3.3: Uncertainty/use: What uncertainty floor and confidence-interval conversion should be used, and how should rows above fixed porosity be interpreted?\\
\addlinespace[0.18em]
\textbf{DQ039} & What is the preferred next NMR residual? & Source inventory 3.3: OGS tie: Should comparison be to OGS saturation S\_l, theta = n*S\_l, corrected/free-water theta, label-bias corrected theta, or within-label trend/anomaly only? Source inventory 3.3: Uncertainty/use: What uncertainty floor and confidence-interval conversion should be used, and how should rows above fixed porosity be interpreted?\\
\addlinespace[0.18em]
\textbf{DQ040} & Which NMR device produced each row? & Source inventory 3.3: Meaning: Does NMR measure total NMR-visible water, mobile/free water, bound/interlayer water, or another hydrogen/water proxy in the supplied CD-A data? Source inventory 3.3: What/how measured: What NMR signal, T2 distribution, calibration, detuning caveat, campaign processing, or label/borehole conversion produces the reported water-content value?\\
\addlinespace[0.18em]
\textbf{DQ041} & Which campaign produced each row? & Source inventory 3.3: Meaning: Does NMR measure total NMR-visible water, mobile/free water, bound/interlayer water, or another hydrogen/water proxy in the supplied CD-A data? Source inventory 3.3: What/how measured: What NMR signal, T2 distribution, calibration, detuning caveat, campaign processing, or label/borehole conversion produces the reported water-content value?\\
\addlinespace[0.18em]
\textbf{DQ042} & What borehole label and depth interval does each row sample? & Source inventory 3.3: Position in domain: Which NMR labels, profiles, borehole points, wall/profile supports, and current Niche 4 mesh supports are active or outside support?\\
\addlinespace[0.18em]
\textbf{DQ043} & What sensitive footprint does one NMR row represent? & Source inventory 3.3: Position in domain: Which NMR labels, profiles, borehole points, wall/profile supports, and current Niche 4 mesh supports are active or outside support?\\
\addlinespace[0.18em]
\textbf{DQ044} & Which fitting model produced the reported water content? & Source inventory 3.3: Meaning: Does NMR measure total NMR-visible water, mobile/free water, bound/interlayer water, or another hydrogen/water proxy in the supplied CD-A data? Source inventory 3.3: What/how measured: What NMR signal, T2 distribution, calibration, detuning caveat, campaign processing, or label/borehole conversion produces the reported water-content value?\\
\addlinespace[0.18em]
\textbf{DQ045} & Which labels are accepted as inside the Niche 4 OGS support? & Source inventory 3.3: Position in domain: Which NMR labels, profiles, borehole points, wall/profile supports, and current Niche 4 mesh supports are active or outside support?\\
\addlinespace[0.18em]
\textbf{DQ046} & Are NMR dates acquisition dates or processed-table dates? & Source inventory 3.3: Time of measurement: Which rows are weekly time series and which are seasonal/campaign profiles, and how do their timestamps map to OGS output times?\\
\addlinespace[0.18em]
\textbf{DQ047} & How should NMR confidence intervals become standard deviations? & Source inventory 3.3: Uncertainty/use: What uncertainty floor and confidence-interval conversion should be used, and how should rows above fixed porosity be interpreted?\\
\addlinespace[0.18em]
\textbf{DQ048} & Does the stored NMR value represent total NMR-visible water? & Source inventory 3.3: Meaning: Does NMR measure total NMR-visible water, mobile/free water, bound/interlayer water, or another hydrogen/water proxy in the supplied CD-A data? Source inventory 3.3: What/how measured: What NMR signal, T2 distribution, calibration, detuning caveat, campaign processing, or label/borehole conversion produces the reported water-content value?\\
\addlinespace[0.18em]
\textbf{DQ049} & Does the stored NMR value represent mobile pore water? & Source inventory 3.3: Meaning: Does NMR measure total NMR-visible water, mobile/free water, bound/interlayer water, or another hydrogen/water proxy in the supplied CD-A data? Source inventory 3.3: What/how measured: What NMR signal, T2 distribution, calibration, detuning caveat, campaign processing, or label/borehole conversion produces the reported water-content value?\\
\addlinespace[0.18em]
\textbf{DQ050} & Can bound water be treated as constant by label? & Source inventory 3.3: Bound water: How should bound/interlayer water not active in OGS saturation be filtered, modelled, estimated, held constant, or represented as bias/uncertainty?\\
\addlinespace[0.18em]
\textbf{DQ051} & Can bound water be treated as constant by campaign or depth? & Source inventory 3.3: Bound water: How should bound/interlayer water not active in OGS saturation be filtered, modelled, estimated, held constant, or represented as bias/uncertainty?\\
\addlinespace[0.18em]
\textbf{DQ052} & Is a T2 threshold available for the mobile fraction? & Source inventory 3.3: Bound water: How should bound/interlayer water not active in OGS saturation be filtered, modelled, estimated, held constant, or represented as bias/uncertainty?\\
\addlinespace[0.18em]
\textbf{DQ053} & Is there a drying, CCM, or other calibration for bound water? & Source inventory 3.3: Bound water: How should bound/interlayer water not active in OGS saturation be filtered, modelled, estimated, held constant, or represented as bias/uncertainty?\\
\addlinespace[0.18em]
\textbf{DQ054} & How should rows with \(\theta_{\mathrm{obs}}>n\) be used? & Source inventory 3.3: Uncertainty/use: What uncertainty floor and confidence-interval conversion should be used, and how should rows above fixed porosity be interpreted?\\
\addlinespace[0.18em]
\textbf{DQ055} & Which forward state does an NMR residual actually depend on? & Source inventory 3.3: OGS tie: Should comparison be to OGS saturation S\_l, theta = n*S\_l, corrected/free-water theta, label-bias corrected theta, or within-label trend/anomaly only? Source inventory 3.3: Uncertainty/use: What uncertainty floor and confidence-interval conversion should be used, and how should rows above fixed porosity be interpreted?\\
\addlinespace[0.18em]
\textbf{DQ056} & What nuisance parameters are needed for absolute NMR use? & Source inventory 3.3: Meaning: Does NMR measure total NMR-visible water, mobile/free water, bound/interlayer water, or another hydrogen/water proxy in the supplied CD-A data? Source inventory 3.3: What/how measured: What NMR signal, T2 distribution, calibration, detuning caveat, campaign processing, or label/borehole conversion produces the reported water-content value?\\
\addlinespace[0.18em]
\textbf{DQ057} & What data-leakage risk must be avoided? & Source inventory 3.3: OGS tie: Should comparison be to OGS saturation S\_l, theta = n*S\_l, corrected/free-water theta, label-bias corrected theta, or within-label trend/anomaly only? Source inventory 3.3: Uncertainty/use: What uncertainty floor and confidence-interval conversion should be used, and how should rows above fixed porosity be interpreted?\\
\addlinespace[0.18em]
\textbf{DQ058} & What does ERT measure? & Source inventory 3.4: Meaning: Is the primary ERT observable resistivity, log resistivity, resistivity change, phase angle, or a derived water-content field? Source inventory 3.4: What/how measured: Is the ERT field an inversion product from electrode data, local apparent resistivity, processed VTK field, or workbook-derived relation, and what correlation structure does the inversion impose? Source inventory 3.4: Units: Are values in ohm m, log10 ohm m, relative resistivity, percent change, or converted water content?\\
\addlinespace[0.18em]
\textbf{DQ059} & How should OGS be tied to ERT? & Source inventory 3.4: Meaning: Is the primary ERT observable resistivity, log resistivity, resistivity change, phase angle, or a derived water-content field? Source inventory 3.4: What/how measured: Is the ERT field an inversion product from electrode data, local apparent resistivity, processed VTK field, or workbook-derived relation, and what correlation structure does the inversion impose? Source inventory 3.4: Units: Are values in ohm m, log10 ohm m, relative resistivity, percent change, or converted water content?\\
\addlinespace[0.18em]
\textbf{DQ060} & Is the ERT residual target resistivity? & Source inventory 3.4: Meaning: Is the primary ERT observable resistivity, log resistivity, resistivity change, phase angle, or a derived water-content field? Source inventory 3.4: What/how measured: Is the ERT field an inversion product from electrode data, local apparent resistivity, processed VTK field, or workbook-derived relation, and what correlation structure does the inversion impose? Source inventory 3.4: Units: Are values in ohm m, log10 ohm m, relative resistivity, percent change, or converted water content?\\
\addlinespace[0.18em]
\textbf{DQ061} & Is the ERT residual target log-resistivity? & Source inventory 3.4: Meaning: Is the primary ERT observable resistivity, log resistivity, resistivity change, phase angle, or a derived water-content field? Source inventory 3.4: What/how measured: Is the ERT field an inversion product from electrode data, local apparent resistivity, processed VTK field, or workbook-derived relation, and what correlation structure does the inversion impose? Source inventory 3.4: Units: Are values in ohm m, log10 ohm m, relative resistivity, percent change, or converted water content?\\
\addlinespace[0.18em]
\textbf{DQ062} & Is the ERT residual target resistivity change? & Source inventory 3.4: Meaning: Is the primary ERT observable resistivity, log resistivity, resistivity change, phase angle, or a derived water-content field? Source inventory 3.4: What/how measured: Is the ERT field an inversion product from electrode data, local apparent resistivity, processed VTK field, or workbook-derived relation, and what correlation structure does the inversion impose? Source inventory 3.4: Units: Are values in ohm m, log10 ohm m, relative resistivity, percent change, or converted water content?\\
\addlinespace[0.18em]
\textbf{DQ063} & Which inversion mesh is authoritative? & Source inventory 3.4: Position in domain: What coordinate transform, projection mesh, near-niche support mask, 35 cm depth/support handling, and aggregation should map ERT cells to the OGS 2D slice?\\
\addlinespace[0.18em]
\textbf{DQ064} & Which VTK field is authoritative? & Source inventory 3.4: Meaning: Is the primary ERT observable resistivity, log resistivity, resistivity change, phase angle, or a derived water-content field? Source inventory 3.4: What/how measured: Is the ERT field an inversion product from electrode data, local apparent resistivity, processed VTK field, or workbook-derived relation, and what correlation structure does the inversion impose? Source inventory 3.4: Units: Are values in ohm m, log10 ohm m, relative resistivity, percent change, or converted water content?\\
\addlinespace[0.18em]
\textbf{DQ065} & Which reference date is used for ERT changes or anomalies? & Source inventory 3.4: Time of measurement: Which ERT timesteps, monthly campaigns, folders, and VTK fields correspond to OGS output times? Source inventory 3.4: Uncertainty/use: What uncertainty/covariance/aggregation model prevents over-weighting dense correlated ERT cells, and should ERT remain diagnostic until transform/support/uncertainty are approved?\\
\addlinespace[0.18em]
\textbf{DQ066} & Are bad electrodes or bad timesteps flagged? & Source inventory 3.4: Time of measurement: Which ERT timesteps, monthly campaigns, folders, and VTK fields correspond to OGS output times? Source inventory 3.4: Uncertainty/use: What uncertainty/covariance/aggregation model prevents over-weighting dense correlated ERT cells, and should ERT remain diagnostic until transform/support/uncertainty are approved?\\
\addlinespace[0.18em]
\textbf{DQ067} & Are inversion artefacts flagged? & Source inventory 3.4: Meaning: Is the primary ERT observable resistivity, log resistivity, resistivity change, phase angle, or a derived water-content field? Source inventory 3.4: What/how measured: Is the ERT field an inversion product from electrode data, local apparent resistivity, processed VTK field, or workbook-derived relation, and what correlation structure does the inversion impose? Source inventory 3.4: Units: Are values in ohm m, log10 ohm m, relative resistivity, percent change, or converted water content?\\
\addlinespace[0.18em]
\textbf{DQ068} & Which water-resistivity law should be used? & Source inventory 3.4: OGS tie: Should OGS saturation or theta = n*S\_l be converted to resistivity through an Archie-type, empirical CD-A-specific, NMR-calibrated, or clay-specific relation? Source inventory 3.4: Bound water and clay conduction: Does bound/locked water conduct current like free pore water, and how do clay surface conduction, salinity, porosity, pore-water conductivity, cracks/faults, and saturation history affect resistivity?\\
\addlinespace[0.18em]
\textbf{DQ069} & What uncertainty is assigned to the water-resistivity exponent? & Source inventory 3.4: OGS tie: Should OGS saturation or theta = n*S\_l be converted to resistivity through an Archie-type, empirical CD-A-specific, NMR-calibrated, or clay-specific relation? Source inventory 3.4: Bound water and clay conduction: Does bound/locked water conduct current like free pore water, and how do clay surface conduction, salinity, porosity, pore-water conductivity, cracks/faults, and saturation history affect resistivity?\\
\addlinespace[0.18em]
\textbf{DQ070} & What pore-fluid resistivity should be used? & Source inventory 3.4: OGS tie: Should OGS saturation or theta = n*S\_l be converted to resistivity through an Archie-type, empirical CD-A-specific, NMR-calibrated, or clay-specific relation? Source inventory 3.4: Bound water and clay conduction: Does bound/locked water conduct current like free pore water, and how do clay surface conduction, salinity, porosity, pore-water conductivity, cracks/faults, and saturation history affect resistivity?\\
\addlinespace[0.18em]
\textbf{DQ071} & What porosity should be used in the ERT transform? & Source inventory 3.4: Meaning: Is the primary ERT observable resistivity, log resistivity, resistivity change, phase angle, or a derived water-content field? Source inventory 3.4: What/how measured: Is the ERT field an inversion product from electrode data, local apparent resistivity, processed VTK field, or workbook-derived relation, and what correlation structure does the inversion impose? Source inventory 3.4: Units: Are values in ohm m, log10 ohm m, relative resistivity, percent change, or converted water content?\\
\addlinespace[0.18em]
\textbf{DQ072} & Does the transform use total water content \(\theta\)? & Source inventory 3.4: OGS tie: Should OGS saturation or theta = n*S\_l be converted to resistivity through an Archie-type, empirical CD-A-specific, NMR-calibrated, or clay-specific relation? Source inventory 3.4: Bound water and clay conduction: Does bound/locked water conduct current like free pore water, and how do clay surface conduction, salinity, porosity, pore-water conductivity, cracks/faults, and saturation history affect resistivity?\\
\addlinespace[0.18em]
\textbf{DQ073} & Does the transform need a free-water correction? & Source inventory 3.4: OGS tie: Should OGS saturation or theta = n*S\_l be converted to resistivity through an Archie-type, empirical CD-A-specific, NMR-calibrated, or clay-specific relation? Source inventory 3.4: Bound water and clay conduction: Does bound/locked water conduct current like free pore water, and how do clay surface conduction, salinity, porosity, pore-water conductivity, cracks/faults, and saturation history affect resistivity?\\
\addlinespace[0.18em]
\textbf{DQ074} & Does clay surface conductivity need an additive term? & Source inventory 3.4: OGS tie: Should OGS saturation or theta = n*S\_l be converted to resistivity through an Archie-type, empirical CD-A-specific, NMR-calibrated, or clay-specific relation? Source inventory 3.4: Bound water and clay conduction: Does bound/locked water conduct current like free pore water, and how do clay surface conduction, salinity, porosity, pore-water conductivity, cracks/faults, and saturation history affect resistivity?\\
\addlinespace[0.18em]
\textbf{DQ075} & Does bound water conduct like mobile pore water? & Source inventory 3.4: OGS tie: Should OGS saturation or theta = n*S\_l be converted to resistivity through an Archie-type, empirical CD-A-specific, NMR-calibrated, or clay-specific relation? Source inventory 3.4: Bound water and clay conduction: Does bound/locked water conduct current like free pore water, and how do clay surface conduction, salinity, porosity, pore-water conductivity, cracks/faults, and saturation history affect resistivity?\\
\addlinespace[0.18em]
\textbf{DQ076} & Can the calibration separate porosity from water content? & Source inventory 3.4: OGS tie: Should OGS saturation or theta = n*S\_l be converted to resistivity through an Archie-type, empirical CD-A-specific, NMR-calibrated, or clay-specific relation? Source inventory 3.4: Bound water and clay conduction: Does bound/locked water conduct current like free pore water, and how do clay surface conduction, salinity, porosity, pore-water conductivity, cracks/faults, and saturation history affect resistivity?\\
\addlinespace[0.18em]
\textbf{DQ077} & Can the calibration separate mineralogy from water content? & Source inventory 3.4: OGS tie: Should OGS saturation or theta = n*S\_l be converted to resistivity through an Archie-type, empirical CD-A-specific, NMR-calibrated, or clay-specific relation? Source inventory 3.4: Bound water and clay conduction: Does bound/locked water conduct current like free pore water, and how do clay surface conduction, salinity, porosity, pore-water conductivity, cracks/faults, and saturation history affect resistivity?\\
\addlinespace[0.18em]
\textbf{DQ078} & Can the calibration separate pore-fluid chemistry from water content? & Source inventory 3.4: OGS tie: Should OGS saturation or theta = n*S\_l be converted to resistivity through an Archie-type, empirical CD-A-specific, NMR-calibrated, or clay-specific relation? Source inventory 3.4: Bound water and clay conduction: Does bound/locked water conduct current like free pore water, and how do clay surface conduction, salinity, porosity, pore-water conductivity, cracks/faults, and saturation history affect resistivity?\\
\addlinespace[0.18em]
\textbf{DQ079} & Is the ERT inversion electrically isotropic? & Source inventory 3.4: Meaning: Is the primary ERT observable resistivity, log resistivity, resistivity change, phase angle, or a derived water-content field? Source inventory 3.4: What/how measured: Is the ERT field an inversion product from electrode data, local apparent resistivity, processed VTK field, or workbook-derived relation, and what correlation structure does the inversion impose? Source inventory 3.4: Units: Are values in ohm m, log10 ohm m, relative resistivity, percent change, or converted water content?\\
\addlinespace[0.18em]
\textbf{DQ080} & What resistivity tensor is used if anisotropy matters? & Source inventory 3.4: Anisotropy: Does electrical anisotropy matter, and should it connect to bedding, fractures/faults, or permeability anisotropy?\\
\addlinespace[0.18em]
\textbf{DQ081} & Does the ERT bedding angle match the hydraulic bedding angle? & Source inventory 3.4: Anisotropy: Does electrical anisotropy matter, and should it connect to bedding, fractures/faults, or permeability anisotropy?\\
\addlinespace[0.18em]
\textbf{DQ082} & Does the ERT bedding angle match the mechanical bedding angle? & Source inventory 3.4: Anisotropy: Does electrical anisotropy matter, and should it connect to bedding, fractures/faults, or permeability anisotropy?\\
\addlinespace[0.18em]
\textbf{DQ083} & Do fractures require a separate ERT mask or term? & Source inventory 3.4: Position in domain: What coordinate transform, projection mesh, near-niche support mask, 35 cm depth/support handling, and aggregation should map ERT cells to the OGS 2D slice?\\
\addlinespace[0.18em]
\textbf{DQ084} & Do carbonate layers require a separate ERT mask or term? & Source inventory 3.4: Position in domain: What coordinate transform, projection mesh, near-niche support mask, 35 cm depth/support handling, and aggregation should map ERT cells to the OGS 2D slice?\\
\addlinespace[0.18em]
\textbf{DQ085} & What ERT-to-OGS coordinate transform is accepted? & Source inventory 3.4: Position in domain: What coordinate transform, projection mesh, near-niche support mask, 35 cm depth/support handling, and aggregation should map ERT cells to the OGS 2D slice?\\
\addlinespace[0.18em]
\textbf{DQ086} & What near-niche ERT mask is accepted? & Source inventory 3.4: Position in domain: What coordinate transform, projection mesh, near-niche support mask, 35 cm depth/support handling, and aggregation should map ERT cells to the OGS 2D slice?\\
\addlinespace[0.18em]
\textbf{DQ087} & Is the 35 cm rock-depth convention part of the ERT support? & Source inventory 3.4: Position in domain: What coordinate transform, projection mesh, near-niche support mask, 35 cm depth/support handling, and aggregation should map ERT cells to the OGS 2D slice?\\
\addlinespace[0.18em]
\textbf{DQ088} & How are dense ERT cells aggregated before scoring? & Source inventory 3.4: Position in domain: What coordinate transform, projection mesh, near-niche support mask, 35 cm depth/support handling, and aggregation should map ERT cells to the OGS 2D slice?\\
\addlinespace[0.18em]
\textbf{DQ089} & What spatial covariance is assigned to ERT cells? & Source inventory 3.4: Position in domain: What coordinate transform, projection mesh, near-niche support mask, 35 cm depth/support handling, and aggregation should map ERT cells to the OGS 2D slice?\\
\addlinespace[0.18em]
\textbf{DQ090} & What temporal covariance is assigned to ERT timesteps? & Source inventory 3.4: Time of measurement: Which ERT timesteps, monthly campaigns, folders, and VTK fields correspond to OGS output times? Source inventory 3.4: Uncertainty/use: What uncertainty/covariance/aggregation model prevents over-weighting dense correlated ERT cells, and should ERT remain diagnostic until transform/support/uncertainty are approved?\\
\addlinespace[0.18em]
\textbf{DQ091} & What is the computational ERT observation operator? & Source inventory 3.4: Meaning: Is the primary ERT observable resistivity, log resistivity, resistivity change, phase angle, or a derived water-content field? Source inventory 3.4: What/how measured: Is the ERT field an inversion product from electrode data, local apparent resistivity, processed VTK field, or workbook-derived relation, and what correlation structure does the inversion impose? Source inventory 3.4: Units: Are values in ohm m, log10 ohm m, relative resistivity, percent change, or converted water content?\\
\addlinespace[0.18em]
\textbf{DQ092} & Why can dense ERT fields dominate an inverse problem incorrectly? & Source inventory 3.4: Uncertainty/use: What uncertainty/covariance/aggregation model prevents over-weighting dense correlated ERT cells, and should ERT remain diagnostic until transform/support/uncertainty are approved?\\
\addlinespace[0.18em]
\textbf{DQ093} & Which ERT uncertainties are mathematically confounded? & Source inventory 3.4: Uncertainty/use: What uncertainty/covariance/aggregation model prevents over-weighting dense correlated ERT cells, and should ERT remain diagnostic until transform/support/uncertainty are approved?\\
\addlinespace[0.18em]
\textbf{DQ094} & What does the Taupe/TDR workbook currently provide? & Source inventory 3.5: Meaning: Does TAUPE mean the Taupe/TDR workbook stream in the repo, and what physical quantity does it represent? Source inventory 3.5: What/how measured: Are values calibrated volumetric water content, apparent relative dielectric permittivity, ARDP/TDR proxy, sensor-specific trend, or qualitative moisture signal? Source inventory 3.5: Units: What are the workbook units, calibration equations, constants, baseline normalization, and sensor/band-specific scaling rules?\\
\addlinespace[0.18em]
\textbf{DQ095} & Which Taupe supports can be mapped to the current OGS slice? & Source inventory 3.5: Position in domain: Which sensors, boreholes, EDZ bands, A3/A4/A7/A8 sheets, line supports, and 2D mesh supports are inside or outside the current model domain?\\
\addlinespace[0.18em]
\textbf{DQ096} & How can Taupe/TDR enter OGS before direct-value calibration is known? & Source inventory 3.5: Units: What are the workbook units, calibration equations, constants, baseline normalization, and sensor/band-specific scaling rules? Source inventory 3.5: OGS tie: Should Taupe/TDR compare to absolute theta = n*S\_l, saturation S\_l, band-averaged theta anomaly, or grouped trend diagnostics only?\\
\addlinespace[0.18em]
\textbf{DQ097} & Does the workbook value store apparent relative dielectric permittivity? & Source inventory 3.5: Units: What are the workbook units, calibration equations, constants, baseline normalization, and sensor/band-specific scaling rules? Source inventory 3.5: OGS tie: Should Taupe/TDR compare to absolute theta = n*S\_l, saturation S\_l, band-averaged theta anomaly, or grouped trend diagnostics only?\\
\addlinespace[0.18em]
\textbf{DQ098} & Does the workbook value store ARDP? & Source inventory 3.5: Meaning: Does TAUPE mean the Taupe/TDR workbook stream in the repo, and what physical quantity does it represent? Source inventory 3.5: What/how measured: Are values calibrated volumetric water content, apparent relative dielectric permittivity, ARDP/TDR proxy, sensor-specific trend, or qualitative moisture signal? Source inventory 3.5: Units: What are the workbook units, calibration equations, constants, baseline normalization, and sensor/band-specific scaling rules?\\
\addlinespace[0.18em]
\textbf{DQ099} & Does the workbook value store calibrated volumetric water content? & Source inventory 3.5: Units: What are the workbook units, calibration equations, constants, baseline normalization, and sensor/band-specific scaling rules? Source inventory 3.5: OGS tie: Should Taupe/TDR compare to absolute theta = n*S\_l, saturation S\_l, band-averaged theta anomaly, or grouped trend diagnostics only?\\
\addlinespace[0.18em]
\textbf{DQ100} & Does the workbook value store travel time or amplitude? & Source inventory 3.5: Time of measurement: What timestamp/date convention and baseline/reference date apply to each Taupe/TDR series and anomaly/trend calculation?\\
\addlinespace[0.18em]
\textbf{DQ101} & What sensor-specific calibration converts Taupe/TDR to water content? & Source inventory 3.5: Position in domain: Which sensors, boreholes, EDZ bands, A3/A4/A7/A8 sheets, line supports, and 2D mesh supports are inside or outside the current model domain?\\
\addlinespace[0.18em]
\textbf{DQ102} & What Taupe/TDR calibration range is valid for Opalinus Clay? & Source inventory 3.5: Units: What are the workbook units, calibration equations, constants, baseline normalization, and sensor/band-specific scaling rules? Source inventory 3.5: OGS tie: Should Taupe/TDR compare to absolute theta = n*S\_l, saturation S\_l, band-averaged theta anomaly, or grouped trend diagnostics only?\\
\addlinespace[0.18em]
\textbf{DQ103} & Which electrode tube does each row represent? & Source inventory 3.5: Position in domain: Which sensors, boreholes, EDZ bands, A3/A4/A7/A8 sheets, line supports, and 2D mesh supports are inside or outside the current model domain?\\
\addlinespace[0.18em]
\textbf{DQ104} & Which EDZ band or depth interval does each row represent? & Source inventory 3.5: Position in domain: Which sensors, boreholes, EDZ bands, A3/A4/A7/A8 sheets, line supports, and 2D mesh supports are inside or outside the current model domain?\\
\addlinespace[0.18em]
\textbf{DQ105} & Are A7 and A8 outside the current OGS mesh? & Source inventory 3.5: Position in domain: Which sensors, boreholes, EDZ bands, A3/A4/A7/A8 sheets, line supports, and 2D mesh supports are inside or outside the current model domain?\\
\addlinespace[0.18em]
\textbf{DQ106} & Are air gaps or tube-contact problems flagged? & Source inventory 3.5: Position in domain: Which sensors, boreholes, EDZ bands, A3/A4/A7/A8 sheets, line supports, and 2D mesh supports are inside or outside the current model domain?\\
\addlinespace[0.18em]
\textbf{DQ107} & Is convergence or settling time flagged for Taupe/TDR? & Source inventory 3.5: Time of measurement: What timestamp/date convention and baseline/reference date apply to each Taupe/TDR series and anomaly/trend calculation?\\
\addlinespace[0.18em]
\textbf{DQ108} & Are failed Taupe/TDR sensor segments flagged? & Source inventory 3.5: Position in domain: Which sensors, boreholes, EDZ bands, A3/A4/A7/A8 sheets, line supports, and 2D mesh supports are inside or outside the current model domain?\\
\addlinespace[0.18em]
\textbf{DQ109} & Does bound water affect Taupe/TDR permittivity like mobile water? & Source inventory 3.5: Units: What are the workbook units, calibration equations, constants, baseline normalization, and sensor/band-specific scaling rules? Source inventory 3.5: OGS tie: Should Taupe/TDR compare to absolute theta = n*S\_l, saturation S\_l, band-averaged theta anomaly, or grouped trend diagnostics only?\\
\addlinespace[0.18em]
\textbf{DQ110} & Is a temperature correction needed for Taupe/TDR? & Source inventory 3.5: Units: What are the workbook units, calibration equations, constants, baseline normalization, and sensor/band-specific scaling rules? Source inventory 3.5: OGS tie: Should Taupe/TDR compare to absolute theta = n*S\_l, saturation S\_l, band-averaged theta anomaly, or grouped trend diagnostics only?\\
\addlinespace[0.18em]
\textbf{DQ111} & Is a salinity or conductivity-loss correction needed for Taupe/TDR? & Source inventory 3.5: Units: What are the workbook units, calibration equations, constants, baseline normalization, and sensor/band-specific scaling rules? Source inventory 3.5: OGS tie: Should Taupe/TDR compare to absolute theta = n*S\_l, saturation S\_l, band-averaged theta anomaly, or grouped trend diagnostics only?\\
\addlinespace[0.18em]
\textbf{DQ112} & Should Taupe/TDR be compared as calibrated direct-value water content? & Source inventory 3.5: Units: What are the workbook units, calibration equations, constants, baseline normalization, and sensor/band-specific scaling rules? Source inventory 3.5: OGS tie: Should Taupe/TDR compare to absolute theta = n*S\_l, saturation S\_l, band-averaged theta anomaly, or grouped trend diagnostics only?\\
\addlinespace[0.18em]
\textbf{DQ113} & Should Taupe/TDR be compared as a standardized anomaly? & Source inventory 3.5: Meaning: Does TAUPE mean the Taupe/TDR workbook stream in the repo, and what physical quantity does it represent? Source inventory 3.5: What/how measured: Are values calibrated volumetric water content, apparent relative dielectric permittivity, ARDP/TDR proxy, sensor-specific trend, or qualitative moisture signal? Source inventory 3.5: Units: What are the workbook units, calibration equations, constants, baseline normalization, and sensor/band-specific scaling rules?\\
\addlinespace[0.18em]
\textbf{DQ114} & What does the current Taupe/TDR residual retain and discard? & Source inventory 3.5: OGS tie: Should Taupe/TDR compare to absolute theta = n*S\_l, saturation S\_l, band-averaged theta anomaly, or grouped trend diagnostics only? Source inventory 3.5: Uncertainty/use: What uncertainty and weighting should be assigned by sensor, band, time, calibration quality, support projection, and grouping?\\
\addlinespace[0.18em]
\textbf{DQ115} & What would be required for calibrated direct-value Taupe/TDR likelihood use? & Source inventory 3.5: OGS tie: Should Taupe/TDR compare to absolute theta = n*S\_l, saturation S\_l, band-averaged theta anomaly, or grouped trend diagnostics only? Source inventory 3.5: Uncertainty/use: What uncertainty and weighting should be assigned by sensor, band, time, calibration quality, support projection, and grouping?\\
\addlinespace[0.18em]
\textbf{DQ116} & What inverse-problem risk comes from using all Taupe bands independently? & Source inventory 3.5: Position in domain: Which sensors, boreholes, EDZ bands, A3/A4/A7/A8 sheets, line supports, and 2D mesh supports are inside or outside the current model domain?\\
\addlinespace[0.18em]
\textbf{DQ117} & What does RH measure and how is it converted? & Source inventory 3.6: What/how measured: How is RH converted through Kelvin equation, and what constants are assumed: temperature, density, gas constant, water molar mass/volume, RH percent/fraction convention, and gas/atmospheric reference? Source inventory 3.6: Units: Are derived quantities RH percent, suction, capillary pressure, gauge liquid pressure, absolute liquid pressure, MPa, or Pa?\\
\addlinespace[0.18em]
\textbf{DQ118} & Does RH link to an OGS input or output? & Source inventory 3.6: What/how measured: How is RH converted through Kelvin equation, and what constants are assumed: temperature, density, gas constant, water molar mass/volume, RH percent/fraction convention, and gas/atmospheric reference? Source inventory 3.6: Units: Are derived quantities RH percent, suction, capillary pressure, gauge liquid pressure, absolute liquid pressure, MPa, or Pa?\\
\addlinespace[0.18em]
\textbf{DQ119} & Which RH sensor defines the open-niche boundary candidate? & Source inventory 3.6: Position in domain: Which RH/suction sensors correspond to the open niche or open twin, and which should inform the niche boundary rather than interior cell residuals?\\
\addlinespace[0.18em]
\textbf{DQ120} & Which temperature sensor should accompany the RH sensor? & Source inventory 3.6: Position in domain: Which RH/suction sensors correspond to the open niche or open twin, and which should inform the niche boundary rather than interior cell residuals?\\
\addlinespace[0.18em]
\textbf{DQ121} & Where is the RH sensor relative to \(\Gamma_E\)? & Source inventory 3.6: Position in domain: Which RH/suction sensors correspond to the open niche or open twin, and which should inform the niche boundary rather than interior cell residuals?\\
\addlinespace[0.18em]
\textbf{DQ122} & Are door openings filtered from the RH data? & Source inventory 3.6: What/how measured: How is RH converted through Kelvin equation, and what constants are assumed: temperature, density, gas constant, water molar mass/volume, RH percent/fraction convention, and gas/atmospheric reference? Source inventory 3.6: Units: Are derived quantities RH percent, suction, capillary pressure, gauge liquid pressure, absolute liquid pressure, MPa, or Pa?\\
\addlinespace[0.18em]
\textbf{DQ123} & Are logger problems filtered from the RH data? & Source inventory 3.6: What/how measured: How is RH converted through Kelvin equation, and what constants are assumed: temperature, density, gas constant, water molar mass/volume, RH percent/fraction convention, and gas/atmospheric reference? Source inventory 3.6: Units: Are derived quantities RH percent, suction, capillary pressure, gauge liquid pressure, absolute liquid pressure, MPa, or Pa?\\
\addlinespace[0.18em]
\textbf{DQ124} & Are high-RH saturation intervals filtered? & Source inventory 3.6: What/how measured: How is RH converted through Kelvin equation, and what constants are assumed: temperature, density, gas constant, water molar mass/volume, RH percent/fraction convention, and gas/atmospheric reference? Source inventory 3.6: Units: Are derived quantities RH percent, suction, capillary pressure, gauge liquid pressure, absolute liquid pressure, MPa, or Pa?\\
\addlinespace[0.18em]
\textbf{DQ125} & Is RH stored as percent or fraction before conversion? & Source inventory 3.6: What/how measured: How is RH converted through Kelvin equation, and what constants are assumed: temperature, density, gas constant, water molar mass/volume, RH percent/fraction convention, and gas/atmospheric reference? Source inventory 3.6: Units: Are derived quantities RH percent, suction, capillary pressure, gauge liquid pressure, absolute liquid pressure, MPa, or Pa?\\
\addlinespace[0.18em]
\textbf{DQ126} & What temperature is used in the Kelvin conversion? & Source inventory 3.6: What/how measured: How is RH converted through Kelvin equation, and what constants are assumed: temperature, density, gas constant, water molar mass/volume, RH percent/fraction convention, and gas/atmospheric reference? Source inventory 3.6: Units: Are derived quantities RH percent, suction, capillary pressure, gauge liquid pressure, absolute liquid pressure, MPa, or Pa?\\
\addlinespace[0.18em]
\textbf{DQ127} & What liquid density is used in the Kelvin conversion? & Source inventory 3.6: What/how measured: How is RH converted through Kelvin equation, and what constants are assumed: temperature, density, gas constant, water molar mass/volume, RH percent/fraction convention, and gas/atmospheric reference? Source inventory 3.6: Units: Are derived quantities RH percent, suction, capillary pressure, gauge liquid pressure, absolute liquid pressure, MPa, or Pa?\\
\addlinespace[0.18em]
\textbf{DQ128} & Does the RH conversion produce capillary pressure? & Source inventory 3.6: What/how measured: How is RH converted through Kelvin equation, and what constants are assumed: temperature, density, gas constant, water molar mass/volume, RH percent/fraction convention, and gas/atmospheric reference? Source inventory 3.6: Units: Are derived quantities RH percent, suction, capillary pressure, gauge liquid pressure, absolute liquid pressure, MPa, or Pa?\\
\addlinespace[0.18em]
\textbf{DQ129} & Does the RH conversion produce gauge liquid pressure? & Source inventory 3.6: What/how measured: How is RH converted through Kelvin equation, and what constants are assumed: temperature, density, gas constant, water molar mass/volume, RH percent/fraction convention, and gas/atmospheric reference? Source inventory 3.6: Units: Are derived quantities RH percent, suction, capillary pressure, gauge liquid pressure, absolute liquid pressure, MPa, or Pa?\\
\addlinespace[0.18em]
\textbf{DQ130} & Does the RH sign convention match \(p_c=-\pl\)? & Source inventory 3.6: What/how measured: How is RH converted through Kelvin equation, and what constants are assumed: temperature, density, gas constant, water molar mass/volume, RH percent/fraction convention, and gas/atmospheric reference? Source inventory 3.6: Units: Are derived quantities RH percent, suction, capillary pressure, gauge liquid pressure, absolute liquid pressure, MPa, or Pa?\\
\addlinespace[0.18em]
\textbf{DQ131} & Was the current \(p_E(t)\) curve derived from RH data? & Source inventory 3.6: What/how measured: How is RH converted through Kelvin equation, and what constants are assumed: temperature, density, gas constant, water molar mass/volume, RH percent/fraction convention, and gas/atmospheric reference? Source inventory 3.6: Units: Are derived quantities RH percent, suction, capillary pressure, gauge liquid pressure, absolute liquid pressure, MPa, or Pa?\\
\addlinespace[0.18em]
\textbf{DQ132} & Was the current \(p_E(t)\) curve manually constructed? & Source inventory 3.6: Curve provenance: Why does the local RH-derived envelope differ from the active seasonal pressure curve, and what source table/script/sensor-screening policy generated the active curve?\\
\addlinespace[0.18em]
\textbf{DQ133} & What smoothing was applied to \(p_E(t)\)? & Source inventory 3.6: Curve provenance: Why does the local RH-derived envelope differ from the active seasonal pressure curve, and what source table/script/sensor-screening policy generated the active curve?\\
\addlinespace[0.18em]
\textbf{DQ134} & What interpolation policy is used for \(p_E(t)\)? & Source inventory 3.6: Time of measurement: What is the RH/T timestamp convention, valid interval, high-RH caution period, low-outlier period, and relation to the active open\_niche\_seasonal curve time axis? Source inventory 3.6: Curve provenance: Why does the local RH-derived envelope differ from the active seasonal pressure curve, and what source table/script/sensor-screening policy generated the active curve?\\
\addlinespace[0.18em]
\textbf{DQ135} & What extrapolation policy is used outside the \(p_E(t)\) date range? & Source inventory 3.6: Time of measurement: What is the RH/T timestamp convention, valid interval, high-RH caution period, low-outlier period, and relation to the active open\_niche\_seasonal curve time axis? Source inventory 3.6: Curve provenance: Why does the local RH-derived envelope differ from the active seasonal pressure curve, and what source table/script/sensor-screening policy generated the active curve?\\
\addlinespace[0.18em]
\textbf{DQ136} & What uncertainty is assigned to \(p_E(t)\)? & Source inventory 3.6: What/how measured: How is RH converted through Kelvin equation, and what constants are assumed: temperature, density, gas constant, water molar mass/volume, RH percent/fraction convention, and gas/atmospheric reference? Source inventory 3.6: Units: Are derived quantities RH percent, suction, capillary pressure, gauge liquid pressure, absolute liquid pressure, MPa, or Pa?\\
\addlinespace[0.18em]
\textbf{DQ137} & Is RH a residual or a forcing parameter in the inverse problem? & Source inventory 3.6: What/how measured: How is RH converted through Kelvin equation, and what constants are assumed: temperature, density, gas constant, water molar mass/volume, RH percent/fraction convention, and gas/atmospheric reference? Source inventory 3.6: Units: Are derived quantities RH percent, suction, capillary pressure, gauge liquid pressure, absolute liquid pressure, MPa, or Pa?\\
\addlinespace[0.18em]
\textbf{DQ138} & Why is RH boundary uncertainty crucial for permeability inversion? & Source inventory 3.6: What/how measured: How is RH converted through Kelvin equation, and what constants are assumed: temperature, density, gas constant, water molar mass/volume, RH percent/fraction convention, and gas/atmospheric reference? Source inventory 3.6: Units: Are derived quantities RH percent, suction, capillary pressure, gauge liquid pressure, absolute liquid pressure, MPa, or Pa?\\
\addlinespace[0.18em]
\textbf{DQ139} & What is the minimum computational representation of \(p_E(t)\)? & Source inventory 3.6: Uncertainty/use: What uncertainty should be assigned if RH becomes a boundary or retention likelihood, and how do we avoid double-counting the same data as both forcing and validation? Source inventory 3.6: OGS tie: Is RH used to reconstruct or check a pressure boundary input, validate active forcing, inform retention parameters, or provide uncertainty on boundary scenarios?\\
\addlinespace[0.18em]
\textbf{DQ140} & What does the direct pressure stream measure? & Source inventory 3.7: Meaning: Do direct pressure, hydraulic-head, pore-pressure, or mini-piezometer measurements exist, including the two more distant sensors remembered from Gesa material? Source inventory 3.7: OGS tie: Can these measurements compare directly to OGS liquid pressure p\_l, or do they require head/elevation/gauge/reference conversion?\\
\addlinespace[0.18em]
\textbf{DQ141} & What are the accepted OGS coordinates for BCD-A28--A31? & Source inventory 3.7: Position in domain: Where are the sensors in coordinates, elevation, borehole/support geometry, distance from the niche, and 2D model frame?\\
\addlinespace[0.18em]
\textbf{DQ142} & What screened interval belongs to each pressure sensor? & Source inventory 3.7: Position in domain: Where are the sensors in coordinates, elevation, borehole/support geometry, distance from the niche, and 2D model frame?\\
\addlinespace[0.18em]
\textbf{DQ143} & What depth or elevation belongs to each pressure sensor? & Source inventory 3.7: Position in domain: Where are the sensors in coordinates, elevation, borehole/support geometry, distance from the niche, and 2D model frame?\\
\addlinespace[0.18em]
\textbf{DQ144} & What exact timestamps are available for the pressure stream? & Source inventory 3.7: Time of measurement: What timestamps, campaigns, maintenance intervals, reference-zero periods, and overlap with simulation outputs exist?\\
\addlinespace[0.18em]
\textbf{DQ145} & What sampling interval is used by the pressure stream? & Source inventory 3.7: Position in domain: Where are the sensors in coordinates, elevation, borehole/support geometry, distance from the niche, and 2D model frame?\\
\addlinespace[0.18em]
\textbf{DQ146} & What maintenance flags exist for the pressure stream? & Source inventory 3.7: Time of measurement: What timestamps, campaigns, maintenance intervals, reference-zero periods, and overlap with simulation outputs exist?\\
\addlinespace[0.18em]
\textbf{DQ147} & Are the exported pressure values pore pressure? & Source inventory 3.7: What/how measured: What pressure quantity is measured: pore pressure, hydraulic head, gauge pressure, absolute pressure, suction/capillary pressure, or a processed Geoscope value? Source inventory 3.7: Units: Are values Pa, MPa, bar, head in m, gauge pressure, absolute pressure, or relative to atmospheric/gas reference?\\
\addlinespace[0.18em]
\textbf{DQ148} & Are the exported pressure values gauge pressure? & Source inventory 3.7: What/how measured: What pressure quantity is measured: pore pressure, hydraulic head, gauge pressure, absolute pressure, suction/capillary pressure, or a processed Geoscope value? Source inventory 3.7: Units: Are values Pa, MPa, bar, head in m, gauge pressure, absolute pressure, or relative to atmospheric/gas reference?\\
\addlinespace[0.18em]
\textbf{DQ149} & Are the exported pressure values hydraulic head? & Source inventory 3.7: What/how measured: What pressure quantity is measured: pore pressure, hydraulic head, gauge pressure, absolute pressure, suction/capillary pressure, or a processed Geoscope value? Source inventory 3.7: Units: Are values Pa, MPa, bar, head in m, gauge pressure, absolute pressure, or relative to atmospheric/gas reference?\\
\addlinespace[0.18em]
\textbf{DQ150} & What pressure unit is used in the raw export? & Source inventory 3.7: What/how measured: What pressure quantity is measured: pore pressure, hydraulic head, gauge pressure, absolute pressure, suction/capillary pressure, or a processed Geoscope value? Source inventory 3.7: Units: Are values Pa, MPa, bar, head in m, gauge pressure, absolute pressure, or relative to atmospheric/gas reference?\\
\addlinespace[0.18em]
\textbf{DQ151} & What atmospheric correction is applied? & Source inventory 3.7: What/how measured: What pressure quantity is measured: pore pressure, hydraulic head, gauge pressure, absolute pressure, suction/capillary pressure, or a processed Geoscope value? Source inventory 3.7: Units: Are values Pa, MPa, bar, head in m, gauge pressure, absolute pressure, or relative to atmospheric/gas reference?\\
\addlinespace[0.18em]
\textbf{DQ152} & What elevation correction is applied? & Source inventory 3.7: Position in domain: Where are the sensors in coordinates, elevation, borehole/support geometry, distance from the niche, and 2D model frame?\\
\addlinespace[0.18em]
\textbf{DQ153} & Which pressure sensors are in the open twin? & Source inventory 3.7: Position in domain: Where are the sensors in coordinates, elevation, borehole/support geometry, distance from the niche, and 2D model frame?\\
\addlinespace[0.18em]
\textbf{DQ154} & Which pressure sensors are in the closed twin? & Source inventory 3.7: Position in domain: Where are the sensors in coordinates, elevation, borehole/support geometry, distance from the niche, and 2D model frame?\\
\addlinespace[0.18em]
\textbf{DQ155} & Which pressure sensors are far-field supports? & Source inventory 3.7: Position in domain: Where are the sensors in coordinates, elevation, borehole/support geometry, distance from the niche, and 2D model frame?\\
\addlinespace[0.18em]
\textbf{DQ156} & Which pressure sensors are outside the 2D slice? & Source inventory 3.7: Position in domain: Where are the sensors in coordinates, elevation, borehole/support geometry, distance from the niche, and 2D model frame?\\
\addlinespace[0.18em]
\textbf{DQ157} & Does the pressure stream validate the far-field Dirichlet pressure? & Source inventory 3.7: Position in domain: Where are the sensors in coordinates, elevation, borehole/support geometry, distance from the niche, and 2D model frame?\\
\addlinespace[0.18em]
\textbf{DQ158} & Does the pressure stream validate the niche boundary evolution? & Source inventory 3.7: OGS tie: Can these measurements compare directly to OGS liquid pressure p\_l, or do they require head/elevation/gauge/reference conversion? Source inventory 3.7: Uncertainty/use: Are they reliable enough for hard residuals, or only qualitative validation because of distance, noise, sensor failures, maintenance, calibration, or 2D-model error?\\
\addlinespace[0.18em]
\textbf{DQ159} & Does the pressure stream validate the initial pressure field? & Source inventory 3.7: OGS tie: Can these measurements compare directly to OGS liquid pressure p\_l, or do they require head/elevation/gauge/reference conversion? Source inventory 3.7: Uncertainty/use: Are they reliable enough for hard residuals, or only qualitative validation because of distance, noise, sensor failures, maintenance, calibration, or 2D-model error?\\
\addlinespace[0.18em]
\textbf{DQ160} & Can the 2D model represent the sampled hydraulic path? & Source inventory 3.7: Position in domain: Where are the sensors in coordinates, elevation, borehole/support geometry, distance from the niche, and 2D model frame?\\
\addlinespace[0.18em]
\textbf{DQ161} & What would the pressure observation operator be? & Source inventory 3.7: OGS tie: Can these measurements compare directly to OGS liquid pressure p\_l, or do they require head/elevation/gauge/reference conversion? Source inventory 3.7: Uncertainty/use: Are they reliable enough for hard residuals, or only qualitative validation because of distance, noise, sensor failures, maintenance, calibration, or 2D-model error?\\
\addlinespace[0.18em]
\textbf{DQ162} & What computational error is most dangerous here? & Source inventory 3.7: OGS tie: Can these measurements compare directly to OGS liquid pressure p\_l, or do they require head/elevation/gauge/reference conversion? Source inventory 3.7: Uncertainty/use: Are they reliable enough for hard residuals, or only qualitative validation because of distance, noise, sensor failures, maintenance, calibration, or 2D-model error?\\
\addlinespace[0.18em]
\textbf{DQ163} & Why are pressure measurements valuable if completed? & Source inventory 3.7: OGS tie: Can these measurements compare directly to OGS liquid pressure p\_l, or do they require head/elevation/gauge/reference conversion? Source inventory 3.7: Uncertainty/use: Are they reliable enough for hard residuals, or only qualitative validation because of distance, noise, sensor failures, maintenance, calibration, or 2D-model error?\\
\addlinespace[0.18em]
\textbf{DQ164} & What does the crackmeter plot measure? & Source inventory 3.9: Meaning: What visible or measurable crack/opening/deformation data exist on the niche surface, separate from large 3D fault geometry? Source inventory 3.9: What/how measured: Are quantities crack width, aperture, crack opening displacement, displacement jump, strain, surface displacement, laser-scan difference, convergence, levelling vertical displacement, pressure response, or qualitative crack class?\\
\addlinespace[0.18em]
\textbf{DQ165} & Is the crackmeter value crack aperture? & Source inventory 3.9: Meaning: What visible or measurable crack/opening/deformation data exist on the niche surface, separate from large 3D fault geometry? Source inventory 3.9: What/how measured: Are quantities crack width, aperture, crack opening displacement, displacement jump, strain, surface displacement, laser-scan difference, convergence, levelling vertical displacement, pressure response, or qualitative crack class?\\
\addlinespace[0.18em]
\textbf{DQ166} & Is the crackmeter value relative displacement? & Source inventory 3.9: Meaning: What visible or measurable crack/opening/deformation data exist on the niche surface, separate from large 3D fault geometry? Source inventory 3.9: What/how measured: Are quantities crack width, aperture, crack opening displacement, displacement jump, strain, surface displacement, laser-scan difference, convergence, levelling vertical displacement, pressure response, or qualitative crack class?\\
\addlinespace[0.18em]
\textbf{DQ167} & Is the crackmeter value strain? & Source inventory 3.9: Meaning: What visible or measurable crack/opening/deformation data exist on the niche surface, separate from large 3D fault geometry? Source inventory 3.9: What/how measured: Are quantities crack width, aperture, crack opening displacement, displacement jump, strain, surface displacement, laser-scan difference, convergence, levelling vertical displacement, pressure response, or qualitative crack class?\\
\addlinespace[0.18em]
\textbf{DQ168} & Does positive value mean opening or closure? & Source inventory 3.9: Units: Are values mm, m, strain, Pa response, relative displacement, survey difference, or qualitative class? Source inventory 3.9: What/how measured: Are quantities crack width, aperture, crack opening displacement, displacement jump, strain, surface displacement, laser-scan difference, convergence, levelling vertical displacement, pressure response, or qualitative crack class?\\
\addlinespace[0.18em]
\textbf{DQ169} & What baseline date is used for the crackmeter trend? & Source inventory 3.9: Meaning: What visible or measurable crack/opening/deformation data exist on the niche surface, separate from large 3D fault geometry? Source inventory 3.9: What/how measured: Are quantities crack width, aperture, crack opening displacement, displacement jump, strain, surface displacement, laser-scan difference, convergence, levelling vertical displacement, pressure response, or qualitative crack class?\\
\addlinespace[0.18em]
\textbf{DQ170} & What instrument axis is measured by each crackmeter? & Source inventory 3.9: Meaning: What visible or measurable crack/opening/deformation data exist on the niche surface, separate from large 3D fault geometry? Source inventory 3.9: What/how measured: Are quantities crack width, aperture, crack opening displacement, displacement jump, strain, surface displacement, laser-scan difference, convergence, levelling vertical displacement, pressure response, or qualitative crack class?\\
\addlinespace[0.18em]
\textbf{DQ171} & Which visible crack does each time series represent? & Source inventory 3.9: Time of measurement: Are data static maps, campaign surveys, continuous Geoscope/crackometer time series, laser-scan epochs, levelling epochs, extensometer series, or reference-zero differences?\\
\addlinespace[0.18em]
\textbf{DQ172} & Should crack data be compared to displacement jump? & Source inventory 3.9: Meaning: What visible or measurable crack/opening/deformation data exist on the niche surface, separate from large 3D fault geometry? Source inventory 3.9: What/how measured: Are quantities crack width, aperture, crack opening displacement, displacement jump, strain, surface displacement, laser-scan difference, convergence, levelling vertical displacement, pressure response, or qualitative crack class?\\
\addlinespace[0.18em]
\textbf{DQ173} & Should crack data be compared to near-niche strain? & Source inventory 3.9: Meaning: What visible or measurable crack/opening/deformation data exist on the niche surface, separate from large 3D fault geometry? Source inventory 3.9: What/how measured: Are quantities crack width, aperture, crack opening displacement, displacement jump, strain, surface displacement, laser-scan difference, convergence, levelling vertical displacement, pressure response, or qualitative crack class?\\
\addlinespace[0.18em]
\textbf{DQ174} & Can visible cracks define a local EDZ prior? & Source inventory 3.9: Material-field tie: Can measured cracks/openings justify local permeability anomaly, porosity/damage zone, hydraulic aperture relation, fracture/EDZ mask, mechanical stiffness reduction, or time-dependent permeability/porosity change?\\
\addlinespace[0.18em]
\textbf{DQ175} & Can visible cracks define a local permeability prior? & Source inventory 3.9: Material-field tie: Can measured cracks/openings justify local permeability anomaly, porosity/damage zone, hydraulic aperture relation, fracture/EDZ mask, mechanical stiffness reduction, or time-dependent permeability/porosity change?\\
\addlinespace[0.18em]
\textbf{DQ176} & What uncertainty should be assigned to crackmeter values? & Source inventory 3.9: Units: Are values mm, m, strain, Pa response, relative displacement, survey difference, or qualitative class? Source inventory 3.9: What/how measured: Are quantities crack width, aperture, crack opening displacement, displacement jump, strain, surface displacement, laser-scan difference, convergence, levelling vertical displacement, pressure response, or qualitative crack class?\\
\addlinespace[0.18em]
\textbf{DQ177} & What OGS quantity could a crackmeter residual compare to? & Source inventory 3.9: OGS tie: Should comparison use OGS displacement, strain, stress, derived crack opening from displacement difference, pressure response, swelling/deformation pattern, or only mechanical plausibility screening? Source inventory 3.9: Uncertainty/use: What numeric exports, support geometry, reference-zero/sign convention, quality flags, registration uncertainty, sensor noise, and 2D projection error are needed before hard residuals?\\
\addlinespace[0.18em]
\textbf{DQ178} & What missing metadata changes the sign of the inverse signal? & Source inventory 3.9: Units: Are values mm, m, strain, Pa response, relative displacement, survey difference, or qualitative class? Source inventory 3.9: What/how measured: Are quantities crack width, aperture, crack opening displacement, displacement jump, strain, surface displacement, laser-scan difference, convergence, levelling vertical displacement, pressure response, or qualitative crack class?\\
\addlinespace[0.18em]
\textbf{DQ179} & How should crack evidence affect permeability candidates now? & Source inventory 3.9: Material-field tie: Can measured cracks/openings justify local permeability anomaly, porosity/damage zone, hydraulic aperture relation, fracture/EDZ mask, mechanical stiffness reduction, or time-dependent permeability/porosity change?\\
\addlinespace[0.18em]
\textbf{DQ180} & What does the precision-levelling summary measure? & Source inventory 3.9: Meaning: What visible or measurable crack/opening/deformation data exist on the niche surface, separate from large 3D fault geometry? Source inventory 3.9: What/how measured: Are quantities crack width, aperture, crack opening displacement, displacement jump, strain, surface displacement, laser-scan difference, convergence, levelling vertical displacement, pressure response, or qualitative crack class?\\
\addlinespace[0.18em]
\textbf{DQ181} & What is still blocked for Geoscope, extensometer, crackmeter, laser, and auxiliary
HM streams? & Source inventory 3.9: OGS tie: Should comparison use OGS displacement, strain, stress, derived crack opening from displacement difference, pressure response, swelling/deformation pattern, or only mechanical plausibility screening? Source inventory 3.9: Uncertainty/use: What numeric exports, support geometry, reference-zero/sign convention, quality flags, registration uncertainty, sensor noise, and 2D projection error are needed before hard residuals?\\
\addlinespace[0.18em]
\textbf{DQ182} & What quantity does each extensometer row measure? & Source inventory 3.9: Meaning: What visible or measurable crack/opening/deformation data exist on the niche surface, separate from large 3D fault geometry? Source inventory 3.9: What/how measured: Are quantities crack width, aperture, crack opening displacement, displacement jump, strain, surface displacement, laser-scan difference, convergence, levelling vertical displacement, pressure response, or qualitative crack class?\\
\addlinespace[0.18em]
\textbf{DQ183} & What quantity does each convergence row measure? & Source inventory 3.9: Meaning: What visible or measurable crack/opening/deformation data exist on the niche surface, separate from large 3D fault geometry? Source inventory 3.9: What/how measured: Are quantities crack width, aperture, crack opening displacement, displacement jump, strain, surface displacement, laser-scan difference, convergence, levelling vertical displacement, pressure response, or qualitative crack class?\\
\addlinespace[0.18em]
\textbf{DQ184} & What quantity does each miniprisma row measure? & Source inventory 3.9: Meaning: What visible or measurable crack/opening/deformation data exist on the niche surface, separate from large 3D fault geometry? Source inventory 3.9: What/how measured: Are quantities crack width, aperture, crack opening displacement, displacement jump, strain, surface displacement, laser-scan difference, convergence, levelling vertical displacement, pressure response, or qualitative crack class?\\
\addlinespace[0.18em]
\textbf{DQ185} & What quantity does each laser-scan product measure? & Source inventory 3.9: Meaning: What visible or measurable crack/opening/deformation data exist on the niche surface, separate from large 3D fault geometry? Source inventory 3.9: What/how measured: Are quantities crack width, aperture, crack opening displacement, displacement jump, strain, surface displacement, laser-scan difference, convergence, levelling vertical displacement, pressure response, or qualitative crack class?\\
\addlinespace[0.18em]
\textbf{DQ186} & What quantity does each levelling row measure? & Source inventory 3.9: Meaning: What visible or measurable crack/opening/deformation data exist on the niche surface, separate from large 3D fault geometry? Source inventory 3.9: What/how measured: Are quantities crack width, aperture, crack opening displacement, displacement jump, strain, surface displacement, laser-scan difference, convergence, levelling vertical displacement, pressure response, or qualitative crack class?\\
\addlinespace[0.18em]
\textbf{DQ187} & What baseline epoch is used by each HM stream? & Source inventory 3.9: Time of measurement: Are data static maps, campaign surveys, continuous Geoscope/crackometer time series, laser-scan epochs, levelling epochs, extensometer series, or reference-zero differences?\\
\addlinespace[0.18em]
\textbf{DQ188} & What sign convention is used by each HM stream? & Source inventory 3.9: Units: Are values mm, m, strain, Pa response, relative displacement, survey difference, or qualitative class? Source inventory 3.9: What/how measured: Are quantities crack width, aperture, crack opening displacement, displacement jump, strain, surface displacement, laser-scan difference, convergence, levelling vertical displacement, pressure response, or qualitative crack class?\\
\addlinespace[0.18em]
\textbf{DQ189} & What coordinate frame is used by each HM stream? & Source inventory 3.9: Meaning: What visible or measurable crack/opening/deformation data exist on the niche surface, separate from large 3D fault geometry? Source inventory 3.9: What/how measured: Are quantities crack width, aperture, crack opening displacement, displacement jump, strain, surface displacement, laser-scan difference, convergence, levelling vertical displacement, pressure response, or qualitative crack class?\\
\addlinespace[0.18em]
\textbf{DQ190} & What point, line, pair, or surface support is measured? & Source inventory 3.9: Position in domain: Where exactly are cracks/openings measured: niche wall location, trace geometry, coordinates, sensor position, open/closed niche side, gauge length, support surface, and 2D slice relation?\\
\addlinespace[0.18em]
\textbf{DQ191} & How is each HM support projected to the 2D slice? & Source inventory 3.9: Position in domain: Where exactly are cracks/openings measured: niche wall location, trace geometry, coordinates, sensor position, open/closed niche side, gauge length, support surface, and 2D slice relation?\\
\addlinespace[0.18em]
\textbf{DQ192} & What instrument covariance is available? & Source inventory 3.9: Meaning: What visible or measurable crack/opening/deformation data exist on the niche surface, separate from large 3D fault geometry? Source inventory 3.9: What/how measured: Are quantities crack width, aperture, crack opening displacement, displacement jump, strain, surface displacement, laser-scan difference, convergence, levelling vertical displacement, pressure response, or qualitative crack class?\\
\addlinespace[0.18em]
\textbf{DQ193} & What survey-network covariance is available? & Source inventory 3.9: Meaning: What visible or measurable crack/opening/deformation data exist on the niche surface, separate from large 3D fault geometry? Source inventory 3.9: What/how measured: Are quantities crack width, aperture, crack opening displacement, displacement jump, strain, surface displacement, laser-scan difference, convergence, levelling vertical displacement, pressure response, or qualitative crack class?\\
\addlinespace[0.18em]
\textbf{DQ194} & How are BCD-A9/B10 failures after September 2025 handled? & Source inventory 3.9: Meaning: What visible or measurable crack/opening/deformation data exist on the niche surface, separate from large 3D fault geometry? Source inventory 3.9: What/how measured: Are quantities crack width, aperture, crack opening displacement, displacement jump, strain, surface displacement, laser-scan difference, convergence, levelling vertical displacement, pressure response, or qualitative crack class?\\
\addlinespace[0.18em]
\textbf{DQ195} & What mechanical observation operators are needed for full HM use? & Source inventory 3.9: OGS tie: Should comparison use OGS displacement, strain, stress, derived crack opening from displacement difference, pressure response, swelling/deformation pattern, or only mechanical plausibility screening? Source inventory 3.9: Uncertainty/use: What numeric exports, support geometry, reference-zero/sign convention, quality flags, registration uncertainty, sensor noise, and 2D projection error are needed before hard residuals?\\
\addlinespace[0.18em]
\textbf{DQ196} & Can these HM streams calibrate permeability in the current frozen workflow? & Source inventory 3.9: Material-field tie: Can measured cracks/openings justify local permeability anomaly, porosity/damage zone, hydraulic aperture relation, fracture/EDZ mask, mechanical stiffness reduction, or time-dependent permeability/porosity change?\\
\addlinespace[0.18em]
\textbf{DQ197} & What would make HM residuals numerically defensible? & Source inventory 3.9: OGS tie: Should comparison use OGS displacement, strain, stress, derived crack opening from displacement difference, pressure response, swelling/deformation pattern, or only mechanical plausibility screening? Source inventory 3.9: Uncertainty/use: What numeric exports, support geometry, reference-zero/sign convention, quality flags, registration uncertainty, sensor noise, and 2D projection error are needed before hard residuals?\\
\addlinespace[0.18em]
\textbf{DQ198} & How should large faults differ from visible niche cracks in the model? & Source inventory 3.8: Position in domain: Which exact .dat file/source contains the 3D fault geometry, coordinate system, units, axes, trace/surface/polygon/point-cloud representation, and relation to the current Niche 4 2D slice? Source inventory 3.8: Units: What units apply to coordinates and aperture/opening if present?\\
\addlinespace[0.18em]
\textbf{DQ199} & Which \texttt{.dat} file defines the large fault geometry? & Source inventory 3.8: Position in domain: Which exact .dat file/source contains the 3D fault geometry, coordinate system, units, axes, trace/surface/polygon/point-cloud representation, and relation to the current Niche 4 2D slice? Source inventory 3.8: Units: What units apply to coordinates and aperture/opening if present?\\
\addlinespace[0.18em]
\textbf{DQ200} & Which Tecplot file defines the \texttt{Kluft} geometry? & Source inventory 3.8: Meaning: Does the faulty/.dat source mean large fault/fault-surface/fracture geometry rather than niche-surface crackmeter data? Source inventory 3.8: Position in domain: Which exact .dat file/source contains the 3D fault geometry, coordinate system, units, axes, trace/surface/polygon/point-cloud representation, and relation to the current Niche 4 2D slice?\\
\addlinespace[0.18em]
\textbf{DQ201} & What coordinate system is used by the fault source file? & Source inventory 3.8: Position in domain: Which exact .dat file/source contains the 3D fault geometry, coordinate system, units, axes, trace/surface/polygon/point-cloud representation, and relation to the current Niche 4 2D slice? Source inventory 3.8: Units: What units apply to coordinates and aperture/opening if present?\\
\addlinespace[0.18em]
\textbf{DQ202} & Does the source contain only geometry? & Source inventory 3.8: Meaning: Does the faulty/.dat source mean large fault/fault-surface/fracture geometry rather than niche-surface crackmeter data? Source inventory 3.8: Position in domain: Which exact .dat file/source contains the 3D fault geometry, coordinate system, units, axes, trace/surface/polygon/point-cloud representation, and relation to the current Niche 4 2D slice?\\
\addlinespace[0.18em]
\textbf{DQ203} & Does the source contain measured aperture? & Source inventory 3.8/2.7: What/how measured: Does the source describe geometry, aperture, hydraulic aperture, mechanical aperture, fracture opening, displacement, fault-zone class, or only visualization geometry? Source inventory 3.8/2.7: Relation to other fields: If aperture is measured, can it be translated to permeability through hydraulic-aperture/cubic-law style relations, and what scale correction and uncertainty are required for clay/fault/EDZ conditions?\\
\addlinespace[0.18em]
\textbf{DQ204} & Does the source contain measured displacement? & Source inventory 3.8: Meaning: Does the faulty/.dat source mean large fault/fault-surface/fracture geometry rather than niche-surface crackmeter data? Source inventory 3.8: Position in domain: Which exact .dat file/source contains the 3D fault geometry, coordinate system, units, axes, trace/surface/polygon/point-cloud representation, and relation to the current Niche 4 2D slice?\\
\addlinespace[0.18em]
\textbf{DQ205} & Does the source contain hydraulic-response information? & Source inventory 3.8/2.7: What/how measured: Does the source describe geometry, aperture, hydraulic aperture, mechanical aperture, fracture opening, displacement, fault-zone class, or only visualization geometry? Source inventory 3.8/2.7: Relation to other fields: If aperture is measured, can it be translated to permeability through hydraulic-aperture/cubic-law style relations, and what scale correction and uncertainty are required for clay/fault/EDZ conditions?\\
\addlinespace[0.18em]
\textbf{DQ206} & Which 3D features intersect the current OGS slice? & Source inventory 3.8: Position in domain: Which exact .dat file/source contains the 3D fault geometry, coordinate system, units, axes, trace/surface/polygon/point-cloud representation, and relation to the current Niche 4 2D slice? Source inventory 3.8: Units: What units apply to coordinates and aperture/opening if present?\\
\addlinespace[0.18em]
\textbf{DQ207} & What effective 2D width is assigned to each fault feature? & Source inventory 3.8: Projection/use: What 3D-to-2D projection is accepted, and how do we avoid overfitting a 2D permeability field to unresolved 3D structures?\\
\addlinespace[0.18em]
\textbf{DQ208} & Should the fault modify permeability? & Source inventory 3.8/2.7: What/how measured: Does the source describe geometry, aperture, hydraulic aperture, mechanical aperture, fracture opening, displacement, fault-zone class, or only visualization geometry? Source inventory 3.8/2.7: Relation to other fields: If aperture is measured, can it be translated to permeability through hydraulic-aperture/cubic-law style relations, and what scale correction and uncertainty are required for clay/fault/EDZ conditions?\\
\addlinespace[0.18em]
\textbf{DQ209} & Should the fault modify porosity? & Source inventory 3.8/2.7: What/how measured: Does the source describe geometry, aperture, hydraulic aperture, mechanical aperture, fracture opening, displacement, fault-zone class, or only visualization geometry? Source inventory 3.8/2.7: Relation to other fields: If aperture is measured, can it be translated to permeability through hydraulic-aperture/cubic-law style relations, and what scale correction and uncertainty are required for clay/fault/EDZ conditions?\\
\addlinespace[0.18em]
\textbf{DQ210} & Should the fault modify EDZ classification? & Source inventory 3.8/2.7: What/how measured: Does the source describe geometry, aperture, hydraulic aperture, mechanical aperture, fracture opening, displacement, fault-zone class, or only visualization geometry? Source inventory 3.8/2.7: Relation to other fields: If aperture is measured, can it be translated to permeability through hydraulic-aperture/cubic-law style relations, and what scale correction and uncertainty are required for clay/fault/EDZ conditions?\\
\addlinespace[0.18em]
\textbf{DQ211} & Should the fault modify tensor orientation? & Source inventory 3.8/2.7: What/how measured: Does the source describe geometry, aperture, hydraulic aperture, mechanical aperture, fracture opening, displacement, fault-zone class, or only visualization geometry? Source inventory 3.8/2.7: Relation to other fields: If aperture is measured, can it be translated to permeability through hydraulic-aperture/cubic-law style relations, and what scale correction and uncertainty are required for clay/fault/EDZ conditions?\\
\addlinespace[0.18em]
\textbf{DQ212} & What uncertainty represents 3D-to-2D projection error? & Source inventory 3.8: Projection/use: What 3D-to-2D projection is accepted, and how do we avoid overfitting a 2D permeability field to unresolved 3D structures?\\
\addlinespace[0.18em]
\textbf{DQ213} & Is the fault static over the experiment? & Source inventory 3.8: Time of measurement: Is the fault dataset static mapped geometry, excavation-time mapping, repeated survey, or time-dependent aperture/activation information?\\
\addlinespace[0.18em]
\textbf{DQ214} & Can the fault open or close mechanically over time? & Source inventory 3.8: Time of measurement: Is the fault dataset static mapped geometry, excavation-time mapping, repeated survey, or time-dependent aperture/activation information?\\
\addlinespace[0.18em]
\textbf{DQ215} & Can the fault self-seal or heal over time? & Source inventory 3.8: Time of measurement: Is the fault dataset static mapped geometry, excavation-time mapping, repeated survey, or time-dependent aperture/activation information?\\
\addlinespace[0.18em]
\textbf{DQ216} & What independent evidence justifies a local permeability anomaly? & Source inventory 3.8/2.7: What/how measured: Does the source describe geometry, aperture, hydraulic aperture, mechanical aperture, fracture opening, displacement, fault-zone class, or only visualization geometry? Source inventory 3.8/2.7: Relation to other fields: If aperture is measured, can it be translated to permeability through hydraulic-aperture/cubic-law style relations, and what scale correction and uncertainty are required for clay/fault/EDZ conditions?\\
\addlinespace[0.18em]
\textbf{DQ217} & How can a fault enter the inverse problem without becoming a fake residual? & Source inventory 3.8: OGS tie: Should the projected fault become a geometry input, structural prior, EDZ/fault material class, permeability or porosity mask, retention/storage prior, mechanical caveat, or qualitative context? Source inventory 3.8: Projection/use: What 3D-to-2D projection is accepted, and how do we avoid overfitting a 2D permeability field to unresolved 3D structures?\\
\addlinespace[0.18em]
\textbf{DQ218} & What is the main identifiability risk of fault-informed inversion? & Source inventory 3.8: Meaning: Does the faulty/.dat source mean large fault/fault-surface/fracture geometry rather than niche-surface crackmeter data? Source inventory 3.8: Position in domain: Which exact .dat file/source contains the 3D fault geometry, coordinate system, units, axes, trace/surface/polygon/point-cloud representation, and relation to the current Niche 4 2D slice?\\
\addlinespace[0.18em]
\textbf{DQ219} & What final decision is needed before using faults quantitatively? & Source inventory 3.8: OGS tie: Should the projected fault become a geometry input, structural prior, EDZ/fault material class, permeability or porosity mask, retention/storage prior, mechanical caveat, or qualitative context? Source inventory 3.8: Projection/use: What 3D-to-2D projection is accepted, and how do we avoid overfitting a 2D permeability field to unresolved 3D structures?\\
\addlinespace[0.18em]
\textbf{DQ220} & Is atmospheric pressure available as a project time series? & Source inventory 3.10/3.11: Other measurements: What additional measurement streams exist in the data beyond those listed here, and for each one what are where, when, units, measured quantity, and OGS input/output/prior/diagnostic status? Source inventory 3.10/3.11: Meaning: Which non-residual support layers are available and required for model use: coordinates/layout, borehole endpoints, mesh projection inputs, bedding/geology, source-to-OGS transforms, and source-file catalogues?\\
\addlinespace[0.18em]
\textbf{DQ221} & Is tunnel temperature available as a project time series? & Source inventory 3.10/3.11: Other measurements: What additional measurement streams exist in the data beyond those listed here, and for each one what are where, when, units, measured quantity, and OGS input/output/prior/diagnostic status? Source inventory 3.10/3.11: Meaning: Which non-residual support layers are available and required for model use: coordinates/layout, borehole endpoints, mesh projection inputs, bedding/geology, source-to-OGS transforms, and source-file catalogues?\\
\addlinespace[0.18em]
\textbf{DQ222} & Is evapometer conductivity available as a project time series? & Source inventory 3.10/3.11: Other measurements: What additional measurement streams exist in the data beyond those listed here, and for each one what are where, when, units, measured quantity, and OGS input/output/prior/diagnostic status? Source inventory 3.10/3.11: Meaning: Which non-residual support layers are available and required for model use: coordinates/layout, borehole endpoints, mesh projection inputs, bedding/geology, source-to-OGS transforms, and source-file catalogues?\\
\addlinespace[0.18em]
\textbf{DQ223} & Are geophone or MSM velocity data available? & Source inventory 3.10/3.11: Other measurements: What additional measurement streams exist in the data beyond those listed here, and for each one what are where, when, units, measured quantity, and OGS input/output/prior/diagnostic status? Source inventory 3.10/3.11: Meaning: Which non-residual support layers are available and required for model use: coordinates/layout, borehole endpoints, mesh projection inputs, bedding/geology, source-to-OGS transforms, and source-file catalogues?\\
\addlinespace[0.18em]
\textbf{DQ224} & Are CCM water-content data available? & Source inventory 3.10/3.11: Other measurements: What additional measurement streams exist in the data beyond those listed here, and for each one what are where, when, units, measured quantity, and OGS input/output/prior/diagnostic status? Source inventory 3.10/3.11: Meaning: Which non-residual support layers are available and required for model use: coordinates/layout, borehole endpoints, mesh projection inputs, bedding/geology, source-to-OGS transforms, and source-file catalogues?\\
\addlinespace[0.18em]
\textbf{DQ225} & Are borehole images available for the current support? & Source inventory 3.10/3.11: Other measurements: What additional measurement streams exist in the data beyond those listed here, and for each one what are where, when, units, measured quantity, and OGS input/output/prior/diagnostic status? Source inventory 3.10/3.11: Meaning: Which non-residual support layers are available and required for model use: coordinates/layout, borehole endpoints, mesh projection inputs, bedding/geology, source-to-OGS transforms, and source-file catalogues?\\
\addlinespace[0.18em]
\textbf{DQ226} & Are geological or lithological logs available for the current support? & Source inventory 3.10/3.11: Other measurements: What additional measurement streams exist in the data beyond those listed here, and for each one what are where, when, units, measured quantity, and OGS input/output/prior/diagnostic status? Source inventory 3.10/3.11: Meaning: Which non-residual support layers are available and required for model use: coordinates/layout, borehole endpoints, mesh projection inputs, bedding/geology, source-to-OGS transforms, and source-file catalogues?\\
\addlinespace[0.18em]
\textbf{DQ227} & Should this stream be an OGS input? & Source inventory 3.10/3.11: OGS tie: Should these streams be treated as observation-operator support, material-field priors, geometry masks, mesh-parameter inputs, or qualitative caveats rather than residuals? Source inventory 3.10/3.11: Activation/use: Which streams are active likelihood terms now, which are diagnostic, which are boundary/input provenance, which are support/prior only, and which are blocked by missing numeric exports or metadata?\\
\addlinespace[0.18em]
\textbf{DQ228} & Should this stream be an OGS output residual? & Source inventory 3.10/3.11: OGS tie: Should these streams be treated as observation-operator support, material-field priors, geometry masks, mesh-parameter inputs, or qualitative caveats rather than residuals? Source inventory 3.10/3.11: Activation/use: Which streams are active likelihood terms now, which are diagnostic, which are boundary/input provenance, which are support/prior only, and which are blocked by missing numeric exports or metadata?\\
\addlinespace[0.18em]
\textbf{DQ229} & Should this stream be a material prior? & Source inventory 3.10/3.11: OGS tie: Should these streams be treated as observation-operator support, material-field priors, geometry masks, mesh-parameter inputs, or qualitative caveats rather than residuals? Source inventory 3.10/3.11: Activation/use: Which streams are active likelihood terms now, which are diagnostic, which are boundary/input provenance, which are support/prior only, and which are blocked by missing numeric exports or metadata?\\
\addlinespace[0.18em]
\textbf{DQ230} & Should this stream be a support mask or quality filter? & Source inventory 3.10/3.11: OGS tie: Should these streams be treated as observation-operator support, material-field priors, geometry masks, mesh-parameter inputs, or qualitative caveats rather than residuals? Source inventory 3.10/3.11: Activation/use: Which streams are active likelihood terms now, which are diagnostic, which are boundary/input provenance, which are support/prior only, and which are blocked by missing numeric exports or metadata?\\
\addlinespace[0.18em]
\textbf{DQ231} & What single missing item blocks use of this stream? & Source inventory 3.10/3.1: Activation/use: Which streams are active likelihood terms now, which are diagnostic, which are boundary/input provenance, which are support/prior only, and which are blocked by missing numeric exports or metadata? Source inventory 3.10/3.1: For each measurement, what preprocessing is required: unit conversion, coordinate transform, 3D-to-2D projection, support averaging, baseline/reference-zero definition, time interpolation, filtering, and uncertainty/correlation model?\\
\addlinespace[0.18em]
\textbf{DQ232} & Who can provide the missing source or decision? & Source inventory 3.10/3.1: Activation/use: Which streams are active likelihood terms now, which are diagnostic, which are boundary/input provenance, which are support/prior only, and which are blocked by missing numeric exports or metadata? Source inventory 3.10/3.1: For each measurement, what preprocessing is required: unit conversion, coordinate transform, 3D-to-2D projection, support averaging, baseline/reference-zero definition, time interpolation, filtering, and uncertainty/correlation model?\\
\addlinespace[0.18em]
\textbf{DQ233} & What triage rule decides whether a newly found stream is usable? & Source inventory 3.10/3.11: OGS tie: Should these streams be treated as observation-operator support, material-field priors, geometry masks, mesh-parameter inputs, or qualitative caveats rather than residuals? Source inventory 3.10/3.11: Activation/use: Which streams are active likelihood terms now, which are diagnostic, which are boundary/input provenance, which are support/prior only, and which are blocked by missing numeric exports or metadata?\\
\addlinespace[0.18em]
\textbf{DQ234} & What should happen if a stream has values but no covariance? & Source inventory 3.10/3.1: Activation/use: Which streams are active likelihood terms now, which are diagnostic, which are boundary/input provenance, which are support/prior only, and which are blocked by missing numeric exports or metadata? Source inventory 3.10/3.1: For each measurement, what preprocessing is required: unit conversion, coordinate transform, 3D-to-2D projection, support averaging, baseline/reference-zero definition, time interpolation, filtering, and uncertainty/correlation model?\\
\addlinespace[0.18em]
\textbf{DQ235} & What source and processing record is required when a stream becomes active? & Source inventory 3.10/3.11: OGS tie: Should these streams be treated as observation-operator support, material-field priors, geometry masks, mesh-parameter inputs, or qualitative caveats rather than residuals? Source inventory 3.10/3.11: Activation/use: Which streams are active likelihood terms now, which are diagnostic, which are boundary/input provenance, which are support/prior only, and which are blocked by missing numeric exports or metadata?\\
\end{longtable}
\footnotesize
% END GENERATED DETAILED MEASUREMENT QUESTION REGISTER

\normalsize
\fi
